
\section{Proofs for overparameterized least-squares}\label{sec proof thm 1}
In this section, we assume the linear Gaussian problem (LGP) of Definition \ref{def LGP}, the overparameterized regime $k=p>n$ and the min-norm model $\bth$ of \eqref{eq:min_norm}. We prove Theorem \ref{thm:master_W2} that derives the asymptotic DC of $\bth$ and we show how this leads to sharp formulae for the risk of the Magnitude- and Hessian-pruned models.


\subsection{Notation and Assumptions}\label{sec:ass_app}
%\noindent\textbf{Notation.}~
For the reader's convenience, we recall some necessary notation and assumptions from Section \ref{sec main}.  We say that a function $f:\R^m\rightarrow\R$ is pseudo-Lipschitz of order $k$, denoted $f\in\rm{PL}(k)$, if there is a constant $L>0$ such that for all $\x,\y\in\R^m$,
$
|f(\x) - f(\y)|\leq L(1+\tn{\x}^{k-1}+\tn{\y}^{k-1})\|\x-\y\|_2 
$
(See also Section \ref{SM useful fact}). 
We say that a sequence of probability distributions $\nu_p$ on $\R^m$ {converges in $W_k$} to 
$\nu$, written $\nu_p\stackrel{W_k}{\Longrightarrow} \nu$, if $W_k(\nu_p,\nu) \rightarrow 0$ as $p \rightarrow \infty$. An equivalent definition is that, for any $f\in\rm{PL}(k)$, $\lim_{p\rightarrow\infty}\E f(X_p)=\E f(X)$, where expectation is with respect to $X_p\sim\nu_p$ and $X\sim\nu$ (e.g., \cite{montanari2017estimation}).  
%
Finally, recall that a sequence of probability distributions $\nu_n$ on $\R^m$ \emph{converges weakly} to $\nu$, if for any bounded Lipschitz function $f$: $\lim_{p\rightarrow\infty}\E f(X_p)=\E f(X)$, where expectation is with respect to $X_p\sim\nu_p$ and $X\sim\nu$.
%
Throughout, we use $C,C',c,c'$ to denote absolute constants (not depending on $n,p$) whose value might change from line to line. 


We focus on a double asymptotic regime where:
$$n,p,s\rightarrow\infty \text{ at fixed overparameterization ratio } \kappa:=p/n>1 \text{ and sparsity level } \alpha:=s/p\in(0,1).$$  
%
For a sequence of random variables $\mathcal{X}_{p}$ that converge in probability to some constant $c$ in the limit of Assumption \ref{ass:linear} below, we  write $\mathcal{X}_{p}\rP c$. For a sequence of event $\Ec_p$ for which $\lim_{p\rightarrow}\Pro(\Ec_p) = 1$, we say that $\Ec_p$ occurs \emph{with probability approaching 1}. For this, we will often use the shorthand ``wpa 1". 

\vspace{5pt}
Next, we recall the set of assumption under which our analysis applies:


\asstwo*
\assthree*


%
%\vp
%\textbf{(A1)[Diagonal covariance].}~~The covariance matrix $\Sigmab$ is diagonal.
%
%%
%\vp
%\textbf{\ref{ass:mu}[Boundedness and empirical distribution].}~~There exist constants $\Sigma_{\min},\Sigma_{\max}\in(0,\infty)$ such that:
%$
%\Sigma_{\min}\leq\bSi_{i,i}\leq \Sigma_{\max}.
%$
%for all $i\in[p].$ 
%%There exists absolute constant  $C>0$ such that $\sqrt{p}|\betas_i|<C$, for all $i\in[p]$. 
%\cts{Furthermore, the joint empirical distribution of $\{(\lai,\sqrt{p}\betas_i)\}_{i\in[p]}$ converges weakly to a probability distribution $\mu$ on $\R_{>0}\times\R$ with bounded $4$th moment 
%%$\E_{(\Lambda,B)}\left[\left(\sqrt{\Lambda^2+B^2}\right)^3\right]$, 
%and assume that as $p\rightarrow\infty$, the $4$th moment of the empirical distribution converges to the $4$th moment of $\mu$, i.e.,
%$$
%\frac{1}{p}\sum_{i=1}^p\left(\sqrt{(\sqrt{p}\betas_i)^{2}+\lai^2}\right)^4 \rP \E_{(\Lambda,B)\sim\mu}\left[\left(\sqrt{\Lambda^2+B^2}\right)^4\right]<\infty.
%$$
%}
% for some $k\geq 3$.}

%\cts{Furthermore, the joint empirical distribution of $\{(\lai,\sqrt{p}\betas_i)\}_{i\in[p]}$ converges weakly to a probability distribution $\mu$ on $\R_{>0}\times\R$ with bounded $(2k-2)$-th moment, and assume that 
%$$
%\frac{1}{p}\sum_{i\in[p]}\sum_{i=1}^p(\sqrt{p}\betas_i)^{2k-2} \rP \E_{(\Lambda,B)\sim\mu}\left[ B^{2k-2}\right]<\infty,
%$$
%as $p\rightarrow\infty$ for some $k\geq 3$.}


\noindent We remark that Assumption \ref{ass:mu} above implies (see \cite[Lem.~4]{bayati2011dynamics} and \cite[Lem.~A3]{javanmard2013state}) that for any pseudo-Lipschitz function $\psi:\R^2\rightarrow\R$ of order $4$, i.e., $\psi\in\rm{PL}(4)$:
$$
\frac{1}{p}\sum_{i=1}^p\psi(\lai,\sqrt{p}\betas_i) \rP \E_{(\Lambda,B)\sim\mu}\left[ \psi(\Lambda,B)\right].
$$




\subsection{Asymptotic distribution and risk characterizations}
%, \fx{which is a special case of Theorem \ref{thm:master_W2}, tailored to our specific need in this paper, that is, characterizing the risk of pruned solutions.} \somm{Theorem 1 should be the special case. Not the other way around!!!}

\fx{In this section, we prove our main result Theorem \ref{thm:master_W2}. Recall that $\bth$ is the min-norm solution. Since the distribution of $\bth$ depends on the problem dimensions (as it is a function of $\X,\y$), when necessary, we will use $\bth_n$ notation to make its dimension dependence explicit.} Let $\bth^P=\Pc(\bth)$ be a pruned version of the min-norm solution $\bth$. Recall from Section \ref{sec:risk}, that the first crucial step in characterizing the risk $\Lc(\bth^P)$ is studying the risk of a threshold-based pruned vector $\Lc(\bth^\Tc_t)$. Theorem \ref{thm:master_W2} below shows how this is possible. 

To keep things slightly more general,  consider $\bth^{g}$ defined such that $\sqrt{p}\bth^{g}=g(\sqrt{p}\bth)$, where $g$ is a Lipschitz function acting entry-wise on $\bth$ (for example, $g$ can be the \fx{(arbitrarily close Lipschitz approximation of the)} thresholding operator $\Tc_t$ of Section \ref{sec:risk}). Then,  the risk of $\bth^g$ can be written as
\begin{align}
\Lc(\bth^g) &= \E_{\Dc}[(\x^T(\bt^\star-\bth^g) + \sigma z)^2] = \sigma^2 + (\bt^\star-\bth^g)^T\Sigmab(\bt^\star-\bth^g) \nn \\
  &= \sigma^2 + \frac{1}{p}\sum_{i=1}^p\Sigmab_{i,i}\big(\sqrt{p}\betas_i-g(\sqrt{p}\bth_i)\big)^2\nn
  \\
    &=: \sigma^2 + \frac{1}{p}\sum_{i=1}^p f\big(\sqrt{p}\bth_i,\sqrt{p}\betas_i,\Sigmab_{i,i}\big),\label{eq:risk_app_f}
% \\
%  &\rP \sigma^2 + \E\left[ \Lambda\left(B-\Tc_t(X_{\kappa,\sigma^2})\right) \right]\,.
\end{align}
where in the last line, we defined function $f:\R^3\rightarrow\R$ as follows:
\begin{align}\label{eq:fdef}
f_\Lc(x,y,z) := \fx{z(y-g(x))}^2.
\end{align}
The following theorem establishes the asymptotic limit of \eqref{eq:risk_app_f}. For the reader's convenience, we repeat the notation introduced in Definition \ref{def:Xi}.
%\begin{definition}[Asymptotic DC -- Overparameterized regime]\label{def:Xi}
Let random variables $(\Lambda,B)\sim \mu$ (where $\mu$ is defined in Assumption \ref{ass:mu}) and fix $\kappa>1$. Define parameter $\xi$ as the unique positive solution to the following equation
%\begin{align}\label{eq:ksi}
$$
\E_{\mu}\Big[ \big({1+(\xi\cdot\Lambda)^{-1}}\big)^{-1} \Big] = {\kappa^{-1}}\,.
$$
Further define the positive parameter $\gamma$ as follows:
$$
\hspace{-0.1in}\gamma := 
%\frac{\sigma^2 + \E_{\mu}\left[\frac{N^2}{(\Lambda^{-1}+\xi)^2}\right]}{1-\kappa\E_{\mu}\left[\frac{1}{\left(1+(\xi\Lambda)^{-1}\right)^2}\right]}  
\Big({\sigma^2 + \E_{\mu}\Big[\frac{B^2\Lambda}{(1+\xi\Lambda)^2}\Big]}\Big)\Big/\Big({1-\kappa\E_{\mu}\Big[\frac{1}{\left(1+(\xi\Lambda)^{-1}\right)^2}\Big]}\Big).
$$
%\ct{@Samet: In your notation $\gamma$ should be $\Gamma\kappa$. Everything matches except the second term in the numerator. In your notation, that second term becomes $\int_0^1\bz^2(t)\bt^2(t) \Sigmab(t) dt$ (note the extra $\Sigmab(t)$ compared to your formula)
%}
With these and $H\sim\Nn(0,1)$, define the random variable
$$
X_{\kappa,\sigma^2}:=X_{\kappa,\sigma^2}(\Lambda,B,H) := \Big(1-\frac{1}{1+ \xi\Lambda}\Big) B + \sqrt{\kappa}\frac{\sqrt{\gamma}\,\Lambda^{-1/2}}{1+(\xi\Lambda)^{-1}} H, 
$$
and let $\Pi_{\kappa,\sigma^2}$ be its distribution.
%\end{definition}


%Recall the definition of the measure $\Pi_{\kappa,\sigma^2}$ in Definition \ref{def:Xi}. Then, 
%$\hat\Pi_n(\y,\X,\betas,\Sigmab)$ converges in Wasserstein-2 distance to $\Pi_{\kappa,\sigma^2}\otimes\mu$.
%Let $\cal{H}$ be the hard-thresholding operation. Draw $\hat{\bt}\sim \Dca(\bt,\bSi)$.
%\[
%\text{risk}_{\bSi^p}(\Hc(\hat{\bt}^p))\rightarrow \text{risk}_{\bSi}(\Hc(\hat{\bt}))
%\]
%\end{theorem}
%
%\begin{corollary} 
%Moreover, for the function $f_\Lc:\R^3\rightarrow\R$, defined in \eqref{eq:fdef} it holds that
%Specifically, for any function $f:\R^k\rightarrow\R$, $f\in\rm{PL}(k)$, it holds that
%where the expectation is over $(\Lambda,B,H)\sim\mu\otimes\Nn(0,1).$
%Specifically, letting $\cal{H}:\R\rightarrow\R$ denote the hard-thresholding operation, it holds that
%\begin{align}
%\| \cal{H}( \hat\w^p )\|_2^2 \rP %\E\left[\left(\cal{H}\left(X_{\kappa,\sigma^2}(\Lambda,N,H) \right)\right)^2\right]
%\end{align}
%%\end{corollary}

\mainthm*

Before we prove the theorem, let us show how it immediately leads to a sharp prediction of the risk behavior. Indeed, a direct application of \eqref{eq:thm} for $f=f_\Lc$ to \eqref{eq:risk_app_f} shows that
\begin{align}
\Lc(\bth^g)\rP \sigma^2 +  \E_{(\Lambda,B,H)\sim\mu\otimes\Nn(0,1)}\left[f_\Lc(X_{\kappa,\sigma^2},B,\Lambda) \right] = \sigma^2 +  \E_{(\Lambda,B,H)\sim\mu\otimes\Nn(0,1)}\left[\Sigma\left(B-g(X_{\kappa,\sigma^2})\right)^2 \right].\label{eq:risk_app_f2}
\end{align}

We further remark on the following two consequences of Theorem \ref{thm:master_W2}. 

First, since \eqref{eq:thm} holds for any $\rm{PL}(2)$ function, we have equivalently that $\hat\Pi_n(\y,\X,\betas,\Sigmab)$ converges in Wasserstein-2 distance to $\Pi_{\kappa,\sigma^2}\otimes\mu$, where recall that $\Pi_{\kappa,\sigma^2}$ is the distribution of the random variable $X_{\kappa,\sigma^2}$.

Second, the theorem implies that the empirical distribution of $\sqrt{p}\bth_n$ converges weakly to $\Pi_{\kappa,\sigma^2}$. To see this, apply \eqref{eq:thm} for the $\rm{PL}(2)$ function $f(x,y,z) = \psi(x)$ where  $\psi:\R\rightarrow\R$ is a bounded Lipschitz test function. 

\subsection{Proof of Theorem \ref{thm:master_W2}}
\vspace{5pt}
Let $\X\in\R^{n\times p}$ have zero-mean and normally distributed rows with a diagonal covariance matrix $\bSi=\E[\x\x^T]$. Given a ground-truth vector $\betas$ and labels $\y=\X\betas+\sigma \z,~\z\sim\Nn(0,\Iden_n)$, we consider the least-squares problem subject to the minimum Euclidian norm constraint (as $\kappa=p/n>1$) given by
\begin{align}\label{eq:PO_beta}
\min_{\bt}\frac{1}{2}\tn{\bt}^2\quad\text{subject to}\quad \y=\X\bt.
\end{align}
%, we solve the least-squares problem  \eqref{optim me2} subject to the . 
It is more convenient to work with the following change of variable: 
\begin{align}\label{eq:w}
\w:=\sqrt{\Sigmab}(\bt-\betas).
\end{align}
 With this, the optimization problem  in \eqref{eq:min_norm} can be rewritten as
\begin{align}\label{eq:PO}
\Phi(\X)=\min_{\w} \frac{1}{2}\tn{\Sigmab^{-1/2}\w+\betas}^2\quad\text{subject to}\quad \Xb\w=\sigma \z,
\end{align}
where we write $\Xb=\X\Sigmab^{-1/2}\distas\Nn(0,1)$. First, using standard arguments, we show that the solution of \eqref{eq:PO} is bounded. Hence, we can constraint the optimization in a sufficiently large compact set without loss of generality.


\begin{lemma}[Boundedness of the solution]\label{lem:bd_PO}
Let $\wh_n:=\wh_n(\X,\z)$ be the minimizer in \eqref{eq:PO}. Then, with probability approaching 1, it holds that $\wh_n\in\Bc$, where
$$\Bc:=\left\{\w\,|\,\|\w\|_2\leq B_{+} \right\},\qquad B_+:=5\sqrt{\frac{\Sigma_{\max}}{\Sigma_{\min}}}\frac{\sqrt{\kappa}+1}{\sqrt{\kappa}-1}(\sqrt{\Sigma_{\max}\E\left[B^2\right]} + \sigma).
% + \sqrt{\Sigma_{\max}\E\left[B^2\right]}.
$$ 
\end{lemma}
\begin{proof}
First, we show that the min-norm solution $\bth=\X^T(\X\X^T)^{-1}\y$ of \eqref{eq:PO_beta} is bounded. Note that $\kappa>1$, thus $\X\X^T$ is invertible wpa 1. We
have,
\begin{align}
\tn{\bth_n}^2 = \y^T(\X\X^T)^{-1}\y \leq \frac{\tn{\y}^2}{\la_{\min}(\X\X^T)}  = \frac{\tn{\y}^2}{\la_{\min}(\Xb\Sigmab\Xb^T)} \leq  \frac{\tn{\y}^2}{\la_{\min}(\Xb\Xb^T)\,\Sigma_{\min}} = \frac{\tn{\y}^2}{\sigma_{\min}^2(\Xb)\,\Sigma_{\min}}. \label{eq:Ubb}
\end{align}
But, wpa 1, 
$
\sigma_{\min}(\Xb)/\sqrt{n} \geq \frac{1}{2}\left(\sqrt{\kappa}-1\right).
$
Furthermore, 
$
\|\y\|_2 \leq \|\Xb\Sigmab^{1/2}\betas\|_2 + \sigma\|\z\|_2 \leq \sigma_{\max}(\Xb)\sqrt{\Sigma_{\max}} \|\betas\|_2 + \sigma\|\z\|_2.
$
Hence, wpa 1,
$$
\|\y\|_2/\sqrt{n} \leq 2(\sqrt{\kappa}+1)\sqrt{\Sigma_{\max}} \sqrt{\E\left[B^2\right]} + 2\sigma,
$$
where we used the facts that wpa 1: $\|z\|_2/\sqrt{n}\rP 1$, $\sigma_{\max}(\Xb)<\sqrt{2n}(\sqrt{\kappa}+1)$ and \cts{by Assumption \ref{ass:mu}}:
$$
\|\betas\|_2^2 = \frac{1}{p}\sum_{i=1}^{p}(\sqrt{p}\betas_i)^2 \rP \E\left[B^2\right].
$$
Put together in \eqref{eq:Ubb}, shows that
\begin{align}
\tn{\bth_n} < \frac{2(\sqrt{\kappa}+1)\sqrt{\Sigma_{\max}} \sqrt{\E\left[B^2\right]} + 2\sigma}{\sqrt{\Sigma_{\min}}(\sqrt{\kappa}-1)/2} =: \tilde{B}_+.\label{eq:bd_beta}
\end{align}
Recalling that $\wh_n= \sqrt{\Sigmab}\bth_n - \sqrt{\Sigmab}\betas$, we conclude, as desired, that wpa 1,
$
\tn{\wh_n} \leq \sqrt{\Sigma_{\max}}\tilde{B}_+ + \sqrt{\Sigma_{\max}} \sqrt{\E[B^2]} \leq B_+.
$
%{\color{red}Missing proof}
\end{proof}

%%%%%%%%%%%%%%%%%%%%%%%%%%%%%%%%%
%%%% Samet's k-th bounded %%%%%%%%%%%%%%%%%
%%\begin{lemma}
%\end{lemma}
\so{\begin{lemma} \label{subexp bth}Let $\bth=\X^\dagger \y$ where $\y=\X\bts+\sigma \z$ where $\X,\z\distas\Nn(0,1)$. As $p,n\rightarrow \infty$, wpa.~1, there exists a constant $C>0$ such that, we have that
\[
\frac{1}{p}\sum_{i=1}^p|\sqrt{p}\bth_i|^k\leq C.
\]
\end{lemma}}
\begin{proof} We first show that entries of $\sqrt{p}\bth$ have finite $k$th moments for any $k\geq 0$. We split the proof in two components by writing $\bth=\ab+\sigma \bb$. 

First, let us focus on $\bb=\X^\dagger\z$. Write singular value decomposition $\X=\Ub\La\V^T$. Then $\bb=\V^T\La\Ub(\Ub\La^2\Ub^T)^{-1}\z=\V^T\La^{-1}\z'$ where $\z'=\Ub^T\z\distas\Nn(0,1)$. First, wpa.~1, $C\sqrt{p} \|\X\|\geq \sigma_{\min}(\X)\gtrsim c\sqrt{p}$ for some constants $C,c>0$ as the covariance $\bSi$ are bounded above and below by constants. To show subgaussianity, we fix unit length $\vb$ and consider the variable
\[
X=\vb^T\bb=\vb^T\V^T\La^{-1}\z'.
\]
by fixing all but $\z'$. Since $\z'$ is Gaussian, the subgaussian norm obeys 
\begin{align}
\tsub{\bb}=\sup_{\tn{\vb}=1}\tsub{X}&\lesssim \sup_{\tn{\vb}=1}\|\vb^T\V^T\La^{-1}\|_{\ell_\infty}\\
&\leq\sup_{\tn{\vb}=1} \tn{\vb^T\V^T}\sigma_{\min}(\X)^{-1}\\
&\leq \sigma_{\min}(\X)^{-1}\\
&\lesssim 1/\sqrt{p}.
\end{align}

To proceed, we address the subgaussianity of the main term 
\begin{align}
\ab&=\X^\dagger\X\bts=\X^T(\X\X^T)^{-1}\X\bts\\
&=\sqrt{\bSi}\Xb^T(\Xb\bSi\Xb^T)^{-1}\Xb\sqrt{\bSi}\bts\\
&=\sqrt{\bSi}\Qb^T (\Qb\bSi\Qb^T)^{-1}\Qb\sqrt{\bSi}\bts\\
&=\Vb^T\Vb\bts\\
&=\sum_{i=1}^n \vb_i\vb_i^T\bts.
\end{align} 
Here, $\Vb$ is the unitary matrix that is the range space of $\X$ and $\Qb$ is the unitary matrix that is the row space of $\Xb$. $(\vb_i)_{i=1}^n$ are rows of $\Vb$. Observe have that the range space of $\Vb$ is statistically identical to the range space of $\Qb\sqrt{\bSi}$. We will work with the properties of the row space $\Vb$ rather than the original design matrix $\X$. 

\noindent\textbf{Expectation of $\ab$:} We claim that $\E[\ab]=\m\odot \bts$ where $1\leq \m_j\geq 0$ for all $1\leq j\leq p$. Thus, by assumption on $\bts$, this would mean that $\E[|\sqrt{p}\E[\ab_i]|^k]$ and $p^{-1}\E[\|\sqrt{p}\E[\ab]\|_{\ell_k}^k]$ is bounded by a constant wpa.~1. To see this, observe that
\[
\E[\ab]=\sum_{i=1}^n\E[\vb_i\vb_i^T]\bts=\sum_{j=1}^p\sum_{i=1}^n\E[\vb_i\vb_i^T]\bts_j \eb_j,
\]
where $\eb_j$ is the $j$'th element of the standard basis. Note that the entries of $\vb_i$ have i.i.d.~Bernoulli signs same as the entries of the rows of $\Qb$. Thus $\E[\vb_i\vb_i^T]$ is a diagonal matrix $\Db^i$. This implies that, $\E[\X^\dagger\X]$ is diagonal and the $j$th index of $\bts$ is the only entry responsible from $\ab_j$. Thus
\[
\m_j=\sum_{i=1}^n \Db^i_{j,j}.
\]
Finally, $\sum_{i=1}^n \Db^i_{j,j}$ is equal to $\tn{\vb^\text{col}_j}^2\leq 1$ where $\vb^\text{col}_j$ is the $j$th column of $\Vb$ proving the result.

\noindent\textbf{Subexponentiality of $\ab$:} To conclude the proof, we show that $\ab$ is subexponential. Fix a unit length vector $\eb$ from standard basis. The claim is same as proving $\te{\sqrt{p}(X-\E[X])}\lesssim 1$ where
\begin{align}
X&=\eb^T\X^\dagger\X\bts=\eb^T\sqrt{\bSi}\Qb^T (\Qb\bSi\Qb^T)^{-1}\Qb\sqrt{\bSi}\bts.
\end{align}
\so{We will focus on the scenario $\bSi=\Iden_p$. General $\bSi$ is handled by the fact that $\bSi$ is upper and lower bounded by identity (thus isometric).} Thus, consider
\begin{align}
X&=\eb^T\X^\dagger\X\bts=\eb^T\Qb^T\Qb\bts\\
&=\frac{\tn{\Qb(\bts+\eb)}^2-\tn{\Qb(\bts-\eb)}^2}{4}.
\end{align}
Thus, our claim is same as showing
\[
\te{\tn{\Qb\vb}^2-\E[\tn{\Qb\vb}^2]}\lesssim 1/\sqrt{p}.
\]
for fixed $\vb=\bts+\vb$. Here, we can swap the randomness and assume fixed $\Qb$ and uniformly random $\vb$ entries of $\vb$ are equal to $\vb_i=\tn{\vb}\frac{\g_i}{\tn{\g}}$ for $\g\sim\Nn(0,\Iden_p)$ and $\tn{\vb}\lesssim 1$. Specifically, fixing $\Qb$ to be first $n$ elements of standard basis, we obtain
\[
\frac{1}{\tn{\vb}^2}(\tn{\Qb\vb}^2-\E[\tn{\Qb\vb}^2])=\frac{\sum_{i=1}^n \g_i^2}{\tn{\g}^2}-\frac{n}{p}=\frac{\sum_{i=1}^n \g_i^2-\frac{n}{p}\tn{\g}^2}{\tn{\g}^2}.
\]
To proceed, observe that both numerator and denominator are sub-exponential variables and concentrate accordingly to find that
\[
\Pro(|\tn{\Qb\vb}^2-\E[\tn{\Qb\vb}^2]|\gtrsim t/\sqrt{p})\leq \e^{-t}.
\]
Thus, entries of $\ab$ obeys $\te{X\sqrt{p}}\lesssim 1$ and together with bounded expectations, we obtain $\E[|\sqrt{p}X|^k]\lesssim 1$.

Finally, combining this with the $\bb$ term, we find the desired result $\E[|\sqrt{p}\bth_i|^k]\lesssim 1$.
%The final line follows from the fact that, since $\bSi_{i,i}$ are bounded below, wpa.~1, $\sigma_{\min}(\X)^{-1}\lesssim 1/\sqrt{p}$.

\so{To conclude with the proof, we use two things. First, $\frac{1}{p}\sum_{i=1}^p\bth_i$ weakly converges to the asymptotic limit given by Def.~\ref{def:Xi}. The proof is omitted for brevity but follows the exact same line of argument as the proof of Theorem \ref{thm:master_W2}. The main difference is that, for weak convergence proof, we don't require $k$th moments or pseudo-Lipschitz conditions (thus, the proof does not require this lemma and there is no circular logic). Once the weak convergence is established, we use Lemma 5 of \cite{bayati2011dynamics}. Since $k$th moments of $\bth_i$ are bounded and $\bth$ is converging weakly, this lemma implies that $\|\bth\|_{\ell_k}^k$ is bounded wpa.~1.}
\end{proof}



\cmt{
\[
\te{\tn{\Vb\eb}^2-\E[\tn{\Vb\eb}^2]}\lesssim 1/\sqrt{p}.
\]
}

\cmt{
Set $\A= (\Qb\bSi\Qb^T)^{-1/2}$. Observe that $\A\lesssim \Iden_n$ and that $\Vb=\A\Qb\sqrt{\bSi}$.
Set $\bar{X}=\eb^T\Qb^T\Qb\bts$. Observing $\Qb\bSi\Qb^T\gtrsim \Sigma_{\min}\Iden_{n}$, this yields that
\[
\E|X|^k\lesssim 
\]
Note that $\E[X]\lesssim 1/\sqrt{p}$ from above. Secondly, }









\cmt{
Let $\Qb$ be a uniformly random subspace satisfying $\Qb\Qb^T=\Iden_n$. Writing $\X=\bar{\X}\sqrt{\bSi}$ and noticing row space of $\bar{\X}$ is statistically identical to $\Qb$,
Now, let
%
\[
\A=\E[\Qb\bSi\Qb^T]^{-1/2}.
\]
Observe that, upper and lower eigenvalues of $\A$ are bounded by constants. Additionally, it follows from the properties of random tight frames that $\A$ is a diagonal matrix (specifically, using rows of $\Qb$ are symmetrically distributed). Now, observe that
\[
\A\Qb\bSi\Qb^T\A=\Iden.
\] 
Consequently $\A\Qb\sqrt{\bSi}\sim \Vb$. Thus, we can write
\[
\Vb^T\Vb\bt=\sqrt{\bSi}\Qb\A^2\Qb^T\sqrt{\bSi}\bt.
\]
We shall argue that, this vector is subexponential. First, given fixed vector $\vb$, we have that
\[
\E[\Qb\A^2\Qb^T\vb]
\]
we bound the infinity norm of the mean given by%$\|\E[\Vb^T\Vb\bt]\|_{\ell_{\infty}}$. Clearly
\[
\|\E[\Vb^T\Vb\bt]\|_{\ell_{\infty}}\leq 
\] }
%%%%%%%%%%%%%%%%%%%%%%%%%%%%%%%%%


Lemma \ref{lem:bd_PO} implies that nothing changes in \eqref{eq:PO} if we further constrain $\w\in\Bc$ in \eqref{eq:PO}. Henceforth, with some abuse of notation, we let
\begin{align}\label{eq:PO_bd}
\Phi(\X):=\min_{\w\in\Bc} \frac{1}{2}\tn{\Sigmab^{-1/2}\w+\betas}^2\quad\text{subject to}\quad \Xb\w=\sigma \z,
\end{align}

Next, in order to analyze the primary optimization (PO) problem in \eqref{eq:PO_bd} in apply the CGMT \cite{thrampoulidis2015lasso}. Specifically, we use the constrained formulation of the CGMT given by Theorem \ref{thm closed}.  Specifically, the auxiliary problem (AO) corresponding to \eqref{eq:PO_bd} takes the following form with $\g\sim\Nn(0,\Iden_n)$, $\h\sim\Nn(0,\Iden_p)$, $h\sim \Nn(0,1)$:% (non-rigorously)
\begin{align}
\phi(\g,\h) = \min_{\w\fx{\in\Bc}} \frac{1}{2}\tn{\Sigmab^{-1/2}\w+\betas}^2\quad\text{subject to}\quad \tn{\g}\tn{\w~{\sigma}}\leq \h^T\w+\sigma h.\label{eq:AO_con}
\end{align}%\tn{\h}
%Set $\bar{\h}=\h/\sqrt{p}$.

We will prove the following techincal result about the AO problem.
\somm{Re-specify the assumptions of this lemma!}
\begin{lemma}[Properties of the AO -- Overparameterized regime]\label{lem:AO}
Let $\phi_n=\phi(\g,\h)$ be the optimal cost of the minimization in \eqref{eq:AO_con}. Define $\bar\phi$ as the optimal cost of the following deterministic min-max problem
\begin{align}\label{eq:AO_det}
\bar\phi:=\max_{u\geq 0}\min_{\tau>0}~ \Dc(u,\tau):=\frac{1}{2}\left({u\tau} + \frac{u\sigma^2}{\tau} - u^2\kappa\,\E\left[ \frac{1}{\Lambda^{-1}+\frac{u}{\tau}} \right] - \fx{\E\left[\frac{B^2}{1+\frac{u}{\tau}\Lambda}\right]} \right).
\end{align}
The following statements are true.

\noindent{(i).}~The AO minimization in \eqref{eq:AO_con} is $\frac{1}{\Sigma_{\max}}$-strongly convex and has a unique minimizer $\wh_n:=\wh_n(\g,\h)$.

\noindent{(ii).}~In the limit of $n,p\rightarrow\infty, p/n=\kappa$, it holds that $\phi(\g,\h)\rP\bar\phi$, i.e., for any $\eps>0$:
$$
\lim_{n\rightarrow\infty}\P\left(|\phi(\g,\h)-\bar\phi|>\eps\right) = 0.
$$

\noindent{(iii).} The max-min optimization in \eqref{eq:AO_det} has a unique saddle point $(u_*,\tau*)$ satisfying the following: 
$$
u_*/\tau_* = \xi\quad\text{and}\quad\tau_* = \gamma,
$$
where $\xi, \gamma$ are defined in Definition \ref{def:Xi}.


\noindent{(iv).}~Let $f:\R^3\rightarrow\R$ be a $\rm{PL}({k})$ function. Let $\bth_n=\Sigmab^{-1/2}\wh_n + \betas$. Then,
$$
\frac{1}{p}\sum_{i=1}^{p}f\left(\sqrt{p}\bth_n(\g,\h),\sqrt{p}\betas,\Sigmab\right) \rP \E_{(B,\Lambda,H)\sim\mu\otimes	\Nn(0,1)}\left[f\left(X_{\kappa,\sigma^2}(B,\Lambda,H),B,\Lambda\right) \right].
$$


\noindent{(v).}~\cts{The empirical distribution of $\bth_n$ converges weakly to the measure of $X_{\kappa,\sigma^2}$, and also,  for large enough absolute constant $C>0$:
\begin{align}\label{eq:k_AO}
\sum_{i=1}^{p}\bth^2_n(\g,\h) < C.
% \rP \E\left[(X_{\kappa,\sigma^2})^{2k-2}\right] < \infty.
\end{align}
}



\end{lemma}




%\begin{lemma}\label{lem:deviation}
%Let $\wh_n=\wh_n(\X,\z)$ be a minimizer of the PO in \eqref{eq:PO}. 
%%Further let $\tau_n=\tau_n(\g,\h,\z)$ and $u_n=u_n(\g,\h,\z)$ be the unique saddle point of the min-max optimization in \eqref{eq:AO_4}. Define $\wh_n=\wh_n(\tau_n,u_n)$ as in \eqref{eq:w_n}. 
%For any pseudo-Lipschitz function $f:\R\rightarrow\R$ of order 2, it hold that
%\begin{align}
%\frac{1}{p} \sum_{i=1}^{p} f\left(\hat\w_{n,i}\right) \rP \E_\mu\left[X_{\kappa,\sigma^2}(\Lambda,N,H)  \right].
%\end{align}
%\end{lemma}


%\begin{proof}[Proof of Lemma \ref{lem:deviation}]
We prove Lemma \ref{lem:AO} in Section \ref{sec:proofAO}. Here, we show how this leads to the proof of Theorem \ref{thm:master_W2} when combined with the CGMT framework \cite{thrampoulidis2015lasso}.

\cts{Let $f:\R^3\rightarrow\R$ be a $\rm{PL}(2)$ function, or, the function $f_\Lc$ defined in \eqref{eq:fdef}.} For convenience, 
 define 
$$F_n(\bth,\betas,\Sigmab):=\frac{1}{p} \sum_{i=1}^{p} f\left(\sqrt{p}\bth_{n,i},\sqrt{p}\betas_{i},\Sigmab_{ii}\right)\quad\text{and}\quad\alpha_*:=\E_\mu\left[f(X_{\kappa,\sigma^2}(\Lambda,B,H),B,\Lambda)\right].$$
Fix  any $\eps>0$ and define the set 
\begin{align}\label{eq:S_set}
\Sc = \Sc(\betas,\Sigmab) =\{\bt\bgl |F_n(\bth_n,\betas,\Sigmab)-\alpha_*|\geq 2\eps\}.
\end{align}
It suffices to prove that the solution $\bth_n$ of the PO in \eqref{eq:PO_beta} satisfies $\bth_n\not\in\Sc$ wpa 1. To see that this is sufficient, note the following. On the one hand, setting $f=f_\Lc$, directly proves \eqref{eq:thm}. On the other hand, recall that $W_2$-convergence is equivalent to convergence of any $\rm{PL}(2)$ test function $f$. 

%\ct{Note that this $\Sc$ is a random set now.}
To prove the desired, we need to consider the ``perturbed" PO and AO problems (compare to \eqref{eq:PO} and \eqref{eq:AO_con}) as:
\begin{align}\label{eq:PO_S}
\Phi_S(\X)=\min_{\w\in\Sc} \frac{1}{2}\tn{\Sigmab^{-1/2}\w+\betas}^2\quad\text{subject to}\quad \Xb\w=\sigma \z,
\end{align}
and
\begin{align}
\phi_S(\g,\h)=\min_{\w\in\Sc} \frac{1}{2}\tn{\Sigmab^{-1/2}\w+\betas}^2\quad\text{subject to}\quad \tn{\g}\tn{\w~{\sigma}}\leq \h^T\w+\sigma h.\label{eq:AO_S}
\end{align}%\tn{\h}
Recall here, that $\Xb=\X\Sigmab^{-1/2}\distas\Nn(0,1)$, $\g\sim\Nn(0,\Iden_n)$, $\h\sim\Nn(0,\Iden_p)$, $h\sim \Nn(0,1)$ and we have used the change of variables $\w:=\sqrt{\Sigmab}(\bt-\betas)$ for convenience. 

 Using \cite[Theorem 6.1(iii)]{thrampoulidis2018precise} it suffices to find costants $\bar\phi, \bar\phi_S$ and $\eta>0$ such that the following three conditions hold:
\begin{enumerate}
\item $\bar\phi_S \geq \bar\phi + 3\eta$,
\item $\phi(\g,\h) \leq \bar\phi + \eta$, with probability approaching 1,
\item $\phi_S(\g,\h) \geq \bar\phi_S - \eta$, with probability approaching 1.
\end{enumerate}
In what follows, we  explicitly find $\bar\phi, \bar\phi_S,\eta$ such that the three conditions above hold.

\vspace{5pt}
\noindent\underline{Satisfying Condition 2}: Recall the deterministic min-max optimization in \eqref{eq:AO_det}. Choose $\bar\phi=\Dc(u_*,\tau_*)$ be the optimal cost of this optimization. From Lemma \ref{lem:AO}(ii), $\phi(\g,\h)\rP\bar\phi$. Thus, for any $\eta>0$, with probability approaching 1:
\begin{align}\label{eq:phi_lim}
\bar\phi + \eta \geq \phi(\g,\h) \geq \bar\phi - \eta.
\end{align}
Clearly then, Condition 2 above holds for any $\eta>0$. 

\vspace{5pt}
\noindent\underline{Satisfying Condition 3}: Next, we will show that the third condition holds for appropriate $\bar\phi$. 
%Let $\tau_n=\tau_n(\g,\h,\z)$ and $u_n=u_n(\g,\h,\z)$ be the unique saddle point of the min-max optimization in \eqref{eq:AO_4}. Define 
Let $\wh_n=\wh_n(\g,\h)$ be the unique minimizer of \eqref{eq:AO_con} as per Lemma \ref{lem:AO}(i), i.e., $\frac{1}{2}\tn{\Sigmab^{-1/2}\wh_n+\bt}^2 = \phi(\g,\h)$. Again from Lemma \ref{lem:AO}, the minimization in \eqref{eq:AO_con} is $1/\Sigma_{\max}$-strongly convex in $\w$. Here, $\Sigma_{\max}$ is the upper bound on the eigenvalues of $\Sigmab$ as per Assumption \ref{ass:mu}. Thus, for any $\tilde\eps>0$ and any feasible $\w$ the following holds (deterministically):
\begin{align}\label{eq:sc}
\frac{1}{2}\tn{\Sigmab^{-1/2}\w+\bt}^2 \geq \phi(\g,\h) + \frac{\tilde\epsilon^2}{2\fx{\Sigma_{\max}}},~\text{provided that}~ \|\w-\wh_n(\g,\h)\|_2 \geq \tilde\epsilon.
\end{align}

Now, we argue that \fx{wpa 1,}
\begin{align}\label{eq:dev_arg}
%~\bt\in\Sc \text{ implies that }
\fx{\text{for all}~\bt\in \Sc~\text{the corresponding}~\w~\text{obeys}}~\|\w-\wh_n(\g,\h)\|_2\geq \tilde\eps,
\end{align}
for an appropriate value of a constant $\tilde\eps>0$ and $\w=\sqrt{\Sigmab}(\bt-\betas)$. Consider any $\bt\in\Sc$. 

\noindent First, by definition in \eqref{eq:S_set},
$$
|F_n(\bt,\betas,\Sigmab)-\alpha_*| \geq 2\eps.
$$
Second, by Lemma \ref{lem:AO}(iv), with probability approaching 1,
$$
|F(\bth_n(\g,\h),\betas,\Sigmab) - \alpha_*| \leq \epsilon.
$$
Third, we will show that wpa 1, there exists universal constant $C>0$ such that
\begin{align}
|F_n(\bth_n(\g,\h),\betas,\Sigmab) - F_n(\bt,\betas,\Sigmab)| \leq C {\|\bth_n(\g,\h) - \bt\|_2}\label{eq:dev2show}.
\end{align}
Before proving \eqref{eq:dev2show}, let us argue how combining the above three displays shows the desired. Indeed, in that case, wpa 1,
\begin{align*}
2\eps &\leq |F_n(\bt,\betas,\Sigmab)-\alpha_*| \leq |F_n(\bth_n(\g,\h),\betas,\Sigmab) - F_n(\bt,\betas,\Sigmab)| + |F_n(\bth_n(\g,\h),\betas,\Sigmab) - \alpha_*| \\
&\leq \epsilon + C \,\|\bt-\bth_n\|_2. \\
&\qquad\Longrightarrow \|\bt-\bth_n\|_2 \geq {\eps}/{C}=:\hat\eps\\
&\qquad\Longrightarrow \|\w-\wh_n\|_2 \geq \hat\eps\sqrt{\Sigma_{\min}}=:\tilde\eps.
\end{align*}
In the last line above, we recalled that $\bt=\Sigmab^{-1/2}\w+\betas$ and $\Sigmab_{i,i}\geq\Sigma_{\min},~i\in[p]$ by Assumption \ref{ass:mu}. This proves \eqref{eq:dev_arg}.

Next, combining \eqref{eq:dev_arg} and \eqref{eq:sc}, we find that wpa 1,
$
\frac{1}{2}\tn{\Sigmab^{-1/2}\w+\bt}^2 \geq \phi(\g,\h) + \frac{\tilde\epsilon^2}{2C},~\text{for all}~ \w\in\Sc.
$
Thus, 
\begin{align}
\phi_S(\g,\h) \geq \phi(\g,\h) + \frac{\tilde\epsilon^2}{2\Sigma_{\max}}.\nn
\end{align}
%\end{proof}
When combined with \eqref{eq:phi_lim}, this shows that
\begin{align}
\phi_S(\g,\h) \geq \bar\phi + \frac{\tilde\epsilon^2}{2\Sigma_{\max}} - \eta.
\end{align}
Thus, choosing $\bar\phi_S = \bar\phi + \frac{\tilde\epsilon^2}{2\Sigma_{\max}}$ proves the Condition 3 above. 


% by the pseudo-Lipschitz property, boundedness of $\bth_n$ and $\bt$ (see Lemma \ref{lem:bd_PO}) the following holds. 

%%%%%%%%%%%%%%%%%%%%%%%%%%%%%%%%%%%
%%%%%%%%%%%%%.  PL(k) argument!!
%%%%%%%%%%%%%%%%%%%%%%%%%%%%%%%%%%%%%
%To simplify notation, let $\ab_i=(\sqrt{p}\bth_{n,i}(\g,\h),\sqrt{p}\betas_i,\Sigmab_{i,i})$ and $\bb_i=(\sqrt{p}\bt_i,\sqrt{p}\betas_i,\Sigmab_{i,i})$, for $i\in[p]$. Then, using $f\in\rm{PL}(k)$,
%\begin{align}
%|F_n(\bth_n(\g,\h),\betas,\Sigmab) - F_n(\bt,\betas,\Sigmab)| &\leq \frac{L}{p}\sum_{i=1}^p (1+\max\{\|\ab_i\|_2^{k-1},\|\bb_i\|_2^{k-1}\})\|\ab_i-\bb_i\|_2\nn \\
%&= \frac{L}{p}\sum_{i=1}^p \left(1+\max\{\|\ab_i\|_2^{k-1},\|\bb_i\|_2^{k-1}\}\right)\cdot|\sqrt{p}\bth_{n,i}(\g,\h) - \sqrt{p}\bt_i|\nn\\
%&\leq L\left(1+\max\{\frac{1}{p}\sum_{i=1}^p \|\ab_i\|_2^{2k-2},\frac{1}{p}\sum_{i=1}^p\|\bb_i\|_2^{2k-2}\} \right)^{1/2}{\|\bth_n(\g,\h) - \bt\|_2}\label{eq:LipC}\,.
%\end{align}
%The last inequality above follows by applying Cauchy-Schwartz.  
%To reach \eqref{eq:dev2show} from the above, it suffices to show that
%\begin{align}
%\max\{\frac{1}{p}\sum_{i=1}^p \|\ab_i\|_2^{2k-2},\frac{1}{p}\sum_{i=1}^p\|\bb_i\|_2^{2k-2}\} \leq C,\label{eq:leqC}
%\end{align}
%for some universal constant $C$ (that may depend on $k$, but not on $n,p$).
%% $\frac{1}{p}\sum_{i=1}^p \|\ab_i\|_2^{2k-2}<\infty$ and $\frac{1}{p}\sum_{i=1}^p \|\bb_i\|_2^{2k-2}<\infty$ as $p\rightarrow\infty$. But, 
%To prove \eqref{eq:leqC}, note that 
% using $a^{\ell_1}b^{\ell_2}c^{\ell_3}\leq a^{\ell}+b^{\ell}+c^\ell,~\ell=\ell_1+\ell_2+\ell_3$, for some constant $C=C(k)>0$ it holds:
%%\ct{Why is this (or something like this true)?}
%%\som{arithmetic - geometric mean inequality?}
%\begin{align}
%\frac{1}{p}\sum_{i=1}^p \|\bb_i\|_2^{2k-2} {\leq} C\left(\frac{1}{p}\sum_{i=1}^p |\sqrt{p}\bt_{i}|^{2k-2} + \frac{1}{p}\sum_{i=1}^p |\sqrt{p}\betas_{i}|^{2k-2} + \frac{1}{p}\sum_{i=1}^p |\Sigmab_{i,i}|^{2k-2} \right)\label{eq:bdmom}\,.
%\end{align}
%\cts{The terms $\frac{1}{p}\sum_{i=1}^p |\sqrt{p}\betas_{i}|^{2k-2} $ and $ \frac{1}{p}\sum_{i=1}^p |\Sigmab_{i,i}|^{2k-2}$ are bounded by Assumption \ref{ass:mu}.}
%%\ct{Need to add assumption that $ \frac{1}{p}\sum_{i=1}^p |\Sigmab_{i,i}|^{2k-2}<\infty$ and $\frac{1}{p}\sum_{i=1}^p (\sqrt{p}|\betas_{i}|^{2k-2} $. OK by just assuming that $B$ is Definition \ref{def:Xi} is bounded.}. 
%{\color{red}Furthermore we additionally know from Lemma \ref{subexp bth} that, 
%\begin{align}
%\frac{1}{p}\sum_{i=1}^p |\sqrt{p}\bt_{i}|^{2k-2} <\infty
%\end{align}
%}
%%\frac{1}{d}\sum_{i=1}^p |\bt_{i}|^{2k-2}\leq \|\bt\|_2^{2k-2}<C'(k)$, where the last inequality follows from feasibility of $\bt$, i.e., $\w=\sqrt{\Sigmab}(\bt-\betas)\in\Bc$; see Lemma \ref{lem:bd_PO} and \eqref{eq:bd_beta}. 
%These together with \eqref{eq:bdmom} show that $\frac{1}{p}\sum_{i=1}^p \|\bb_i\|_2^{2k-2}<C$ wpa 1. Similarly, we can argue for $\frac{1}{p}\sum_{i=1}^p \|\ab_i\|_2^{2k-2}$, completing the proof of \eqref{eq:leqC}.
%%%%%%%%%%%%%%%%%%%%%%%%%%
%%%%%%%%%%%%%%%%%%%%%%%%%%%%%%%%

\vp
\noindent\fx{\textbf{Perturbation analysis via Pseudo-Lipschitzness:}} To complete the proof, let us now show \eqref{eq:dev2show}. Henceforth, $C$ is used to denote a universal constant whose value can change from line to line.

 First, consider the case $f=f_\Lc$, i.e., $f(x,y,z) = y(x-g(z))^2$. We have the following chain of inequalities:
%that we care about for evaluating the risk (see Lemma 5) we can write \eqref{eq:LipC} simpler as follows:
\begin{align}
|F_n(\bth_n(\g,\h),\betas,\Sigmab) - F_n(\bt,\betas,\Sigmab)| &=
 \frac{1}{p}\sum_{i=1}^p|\lai|\left| (\sqrt{p}\betas_i-g(\sqrt{p}\bth_{n,i}))^2-(\sqrt{p}\betas_i-g(\sqrt{p}\bt_{i}))^2 \right|\nn\\
 &\leq
\Sigma_{\max} \frac{1}{p}\sum_{i=1}^p \left| (\sqrt{p}\betas_i-g(\sqrt{p}\bth_{n,i}))^2-(\sqrt{p}\betas_i-g(\sqrt{p}\bt_{i}))^2 \right|\nn\\
 &\leq
\fx{C} \Sigma_{\max} \frac{1}{p}\sum_{i=1}^p (1+ \|\sqrt{p}[\betas_i,\bth_{n,i}]\|_2 +  \|\sqrt{p}[\betas_i,\bt_{i}]\|_2) \sqrt{p}|\bth_{n,i}-\bt_{i}|\nn\\
 &\leq
\fx{C} \Sigma_{\max}  \Big(1+ \frac{1}{\sqrt{p}}\big(\sum_{i=1}^{p}\|\sqrt{p}[\betas_i,\bth_{n,i}]\|_2^2\big)^{1/2} + \frac{1}{\sqrt{p}}\big(\sum_{i=1}^p\|\sqrt{p}[\betas_i,\bt_{i}]\|_2^2\big)^{1/2}\Big) \|\bth_n-\bt\|_2\nn\\
 &\leq
%C  \left(1+ \|\betas\|_2^2+ \|\bth_n\|_2^2 + \|\bt\|_2^2\right) \|\bth_n-\bt\|_2.
C  \left(1+ \max\{\|\betas\|_2^2,\|\bth_n\|_2^2,\|\bt\|_2^2\}^{1/2} \right) \|\bth_n-\bt\|_2.\label{eq:fcase1}
 \end{align}
 In the second line above, we used boundedness of $\lai$ as per Assumption \ref{ass:mu}. \somm{Consider explaining the PL(2) sentence.}\fx{In the third line, we used the fact that the function $\psi(a,b) = (a-g(b))^2$ is $\rm{PL}(2)$.} The fourth line follows by Cauchy-Schwartz  inequality. Finally, in the last line, we used the elementary fact that $a+b+c\leq 3\max\{a,b,c\}$ for $a=2\sum_{i=1}^p(\betas_i)^2$ and $b=\sum_{i=1}^p\bth_{n,i}^2$ and $c=\sum_{i=1}^p\bt_{i}^2$.
 
Second, consider the case $f\in\rm{PL}(2)$. Let $\ab_i=(\sqrt{p}\bth_{n,i}(\g,\h),\sqrt{p}\betas_i,\Sigmab_{i,i})$ and $\bb_i=(\sqrt{p}\bt_i,\sqrt{p}\betas_i,\Sigmab_{i,i})$, for $i\in[p]$. A similar chain of inequality holds:
\somm{The solution for PL(k) is showing that around $\hat{\bt}$, there is an arbitarily close $\hat{\bt}_{bdd}$ with bounded entries.}
\begin{align}
|F_n(\bth_n(\g,\h),\betas,\Sigmab) - F_n(\bt,\betas,\Sigmab)| &\leq \frac{\fx{C}}{p}\sum_{i=1}^p (1+\max\{\|\ab_i\|_2,\|\bb_i\|_2\})\|\ab_i-\bb_i\|_2\nn \\
&= \frac{\fx{C}}{p}\sum_{i=1}^p \left(1+\max\{\|\ab_i\|_2,\|\bb_i\|_2\}\right)\cdot|\sqrt{p}\bth_{n,i}(\g,\h) - \sqrt{p}\bt_i|\nn\\
&\leq \fx{C}\left(1+\max\{\frac{1}{p}\sum_{i=1}^p \|\ab_i\|_2^{2},\frac{1}{p}\sum_{i=1}^p\|\bb_i\|_2^{2}\} \right)^{1/2}{\|\bth_n(\g,\h) - \bt\|_2}\nn\\
&\leq C\left(1+\max\{\|\bth_n\|_2^2,\|\betas\|_2^2,\|\bt\|_2^2,\frac{1}{p}\sum_{i=1}^p\lai^2\} \right)^{1/2}{\|\bth_n(\g,\h) - \bt\|_2}\nn\\
&\leq C\Sigma_{\max}^2\left(1+\max\{\|\bth_n\|_2^2,\|\betas\|_2^2,\|\bt\|_2^2\} \right)^{1/2}{\|\bth_n(\g,\h) - \bt\|_2}.\label{eq:fcase2}
\end{align}
The first inequality uses the fact that $f\in\rm{PL}(2)$. The second inequality in the third line follows by Cauchy-Schwartz. The last inequality used Assumption \ref{ass:mu} on boundedness of $\lai$.


It follows from either \eqref{eq:fcase1} or \eqref{eq:fcase2} that in order to prove \eqref{eq:dev2show}, we need to show boundedness of the following terms: $\|\bth_n\|_2$, $\|\betas\|_2$ and $\|\bt\|_2$. By feasibility of $\bth_n$ and $\bt$, we know that $\bth_n,\bt\in\Bc$. Thus, the desired $\|\bt\|_2<\infty$ and $\|\bth_n\|_2<\infty$ follow directly by Lemma \ref{lem:bd_PO}. (Alternatively, for $\bth_n$ we conclude the desired by directly applying Lemma \ref{lem:AO}(v)). Finally, to prove $\|\betas\|_2<\infty$, note that
$$
\|\betas\|_2^2 =\frac{1}{p}\sum_{i=1}^p(\sqrt{p}\betas_i)^2,
$$
\cts{which is bounded wpa 1 by Assumption \ref{ass:mu}, which implies bounded second moments of $\sqrt{p}\betas$. } This completes the proof of \eqref{eq:dev2show}, as desired.
%And everything is easy now since $\|\betas\|_2^2$ is bounded as long as second moments of empirical distribution as bounded, $\|\bth_n\|_2^2$ is bounded if $\|\betas\|_2^2$ is bounded, and, $\|\bt\|_2^2$ is bounded, since $\bt\in\Bc$ by Lemma 1. To prove 


\vspace{5pt}
\noindent\underline{Satisfying Condition 1:} To prove Condition 1, we simply pick $\eta$ to satisfy the following
\begin{align}
\bar\phi_S > \bar\phi + 3 \eta ~\Leftarrow~ \bar\phi + \frac{\tilde\epsilon^2}{2\Sigma_{\max}} - \eta \geq \bar\phi + 3 \eta ~\Leftarrow~ \eta \leq \frac{\tilde\epsilon^2}{8\Sigma_{\max}}.\nn
\end{align}


This completes the proof of Theorem \ref{thm:master_W2}.

\subsection{Proof of Lemma \ref{lem:AO}}\label{sec:proofAO}

~~~~
~~~~

\vp
%\noindent\underline{Proof of (i):}~

\subsubsection{Proof of (i).} Strong convexity of the objective function in \eqref{eq:PO} is easily verified by the second derivative test. Note here that we use Assumption \ref{ass:mu} that $\lai\leq\Sigma_{\max},~i\in[p].$ Uniqueness of the solution follows directly from strong convexity. \ct{Strictly speaking we might need to also argue existence, i.e., feasibility of the AO. An indirect way is to show feasibility using the CGMT, but it  seems unnecessarily complicated?}



\vspace{5pt}
\subsubsection{Proof of (ii).}Using Lagrangian formulation, the solution $\wh_n$ to \eqref{eq:AO_con} is the same as the solution to the following:
\begin{align}
\left(\wh_n,u_n\right) :=\arg\min_{\w\in\Bc}\max_{u\geq 0} ~\frac{1}{2}\tn{\Sigmab^{-1/2}\w+\betas}^2 + u \left( \sqrt{\tn{\w}^2+\sigma^2} \tn{\gba} - \sqrt{\kappa}\,{\hba}^T\w + \frac{\sigma h}{\sqrt{n}} \right)\label{eq:AO_2}
\end{align}
where we have: (i) set $\gba := \g/\sqrt{n}$ and $\hba:= \h/\sqrt{p}$; (ii) recalled that $p/n=\kappa$; and, (iii) used $\left(\wh_n,u_n\right)$ to denote the optimal solutions in \eqref{eq:AO_2}. The subscript $n$ emphasizes the dependence of $\left(\wh_n,u_n\right)$ on the problem dimensions. Also note that (even though not explicit in the notation) $\left(\wh_n,u_n\right)$ are random variables depending on the realizations of $\gba,\hba$ and $h$.

 Notice that the objective function above is convex in $\w$ and linear (thus, concave) in $u$. Thus, strong duality holds and we can flip the order of min-max. Moreover, in order to make the objective easy to optimize with respect to $\w$, we use the following variational expression for the square-root term $\sqrt{\tn{\w}^2+\sigma^2}$:
$$
\tn{\gba}\sqrt{\tn{\w}^2+\sigma^2} = \tn{\gba}\cdot\min_{\tau\in[\sigma,\sqrt{\sigma^2+B_+^2}]} \left\{ \frac{\tau}{2} + \frac{\tn{\w}^2+\sigma^2}{2\tau} \right\} = \min_{\tau\in[\sigma,\sqrt{\sigma^2+B_+^2}]} \left\{ \frac{\tau\tn{\gba}^2}{2} + \frac{\sigma^2}{2\tau} + \frac{\tn{\w}^2}{2\tau} \right\},
$$
where $B_+$ is defined in Lemma \ref{lem:bd_PO}. For convenience define the constraint set for the variable $\tau$ as $\Tc':=[\sigma,\sqrt{\sigma^2+B_+^2}]$. For reasons to be made clear later in the proof (see proof of statement (iii)), we consider the (possibly larger) set:
\[
\Tc:=[\sigma,\max\{\sqrt{\sigma^2+B_+^2},2\tau_*\}]\,
\]
where $\tau_*$ is as in the statement of the lemma.

The above lead to the following equivalent formulation of \eqref{eq:AO_2}:
\begin{align}
\left(\wh_n,u_n,\tau_n\right) = \max_{u\geq 0}\min_{\w\in\Bc,\tau\in\Tc} ~ \frac{u\tau\tn{\gba}^2}{2} + \frac{u\sigma^2}{2\tau} + \frac{u\sigma h}{\sqrt{n}} + \min_{\w\in\Bc} \left\{ \frac{1}{2}\tn{\Sigmab^{-1/2}\w+\betas}^2 + \frac{u}{2\tau}\tn{\w}^2 - u \sqrt{\kappa}\,{\hba}^T\w \right\} .\label{eq:AO_3}
\end{align}
%\ct{To be fully rigorous, need to show here that the unconstrained min over $\w$ is the same as the constrained $\w\in\Bc$.}
The minimization over $\w$ is easy as it involves a strongly convex quadratic function. First, note that unconstrained optimal $\w':=\w'(\tau,u)$ (for fixed $(\tau,u)$) is given by
\begin{align}\label{eq:w'}
\w':=\w'(\tau,u) = -\left(\Sigmab^{-1}+\frac{u}{\tau}\Iden\right)^{-1}\left(\Sigmab^{-1/2}\betas-u\sqrt{\kappa}\hba\right), 
\end{align}
and \eqref{eq:AO_3} simplifies to
\begin{align}
\left(u_n,\tau_n\right)=\max_{u\geq 0}\min_{\tau\in\Tc} ~ \frac{u\tau\tn{\gba}^2}{2} + \frac{u\sigma^2}{2\tau} + \frac{u\sigma h}{\sqrt{n}} - \frac{1}{2} \left(\Sigmab^{-1/2}\betas-u\sqrt{\kappa}\hba\right)^T\left(\Sigmab^{-1}+\frac{u}{\tau}\Iden\right)^{-1} \left(\Sigmab^{-1/2}\betas-u\sqrt{\kappa}\hba\right)\,=:\Rc(u,\tau) .\label{eq:AO_4}
\end{align}
\somm{Does $u=0$ scenario create a problem in saddle point uniqueness?} 
It can be checked by direct differentiation and the second-derivative test that the objective function in \eqref{eq:AO_4} is strictly convex in $\tau$ and strictly concave in $u$ over the domain $\{(u,\tau)\in\R_+\times\R_+\}$ \footnote{\fx{To analyze the matrix-vector product term in \eqref{eq:AO_4} for $(\tau,u)$ one can use the fact that $\bSi$ is diagonal. This way, as a function of $u$ and $\tau$ the analysis reduces to the properties of relatively simple functions. For instance, for $\tau$, this function is in the form $f(\tau)=-(a+b/\tau)^{-1}$ for $a,b>0$ which is strictly convex.}}. Thus, the saddle point $(u_n,\tau_n)$ is unique. Specifically, this implies that the optimal $\wh_n$ in \eqref{eq:AO_3} is given by (cf. \eqref{eq:w'})
\begin{align}\label{eq:w_n}
\wh_n=\w'(\tau_n,u_n) = -\left(\Sigmab^{-1}+\frac{u_n}{\tau_n}\Iden\right)^{-1}\left(\Sigmab^{-1/2}\betas-u_n\sqrt{\kappa}\hba\right).
\end{align}
In Lemma \ref{lem:AO}(v) we will prove that wpa 1, in the limit of $p\rightarrow\infty$, $\|\wh_n\|_2\leq C$ for sufficiently large absolute constant $C>0$. Thus, by choosing the upper bound in the definition of $\Bc$ in Lemma \ref{lem:bd_PO} strictly larger than C, guarantees that the unconstrained $\wh_n$ in \eqref{eq:w_n} is feasible in \eqref{eq:AO_3}.

\noindent\fx{\textbf{Asymptotic limit of the key quantities $\tau_n,u_n$:}} In what follows, we characterize the high-dimensional limit of the optimal pair $(u_n,\tau_n)$ in the limit $n,p\rightarrow\infty,~p/n\rightarrow\kappa$.
%Now that we have reached a min-max problem over only scalar variables, we are ready to study its convergence properties in the limit $n,p\rightarrow\infty,~p/n\rightarrow\kappa$. 
We start by analyzing the (point-wise) convergence of $\Rc(u,\tau)$. 
For the first three summands in \eqref{eq:AO_4}, we easily find that
$$
\left\{\frac{u\tau\tn{\gba}^2}{2} + \frac{u\sigma^2}{2\tau} + \frac{u\sigma h}{\sqrt{n}}\right\}~ \rP~\left\{
\frac{u\tau}{2} + \frac{u\sigma^2}{2\tau} \right\}.
$$
Next, we study the fourth summand. First, note that
\begin{align}
(u\sqrt{\kappa}\hba)^T\left(\Sigmab^{-1}+\frac{u}{\tau}\Iden\right)^{-1}(u\sqrt{\kappa}\hba) &= u^2\kappa\,\frac{1}{p}\h^T\left(\Sigmab^{-1}+\frac{u}{\tau}\Iden\right)^{-1}\h \nn\\
&= u^2\kappa\,\frac{1}{p}\sum_{i=1}^{p}\frac{\h_i^2}{\lai^{-1}+\frac{u}{\tau}} \nn\\
&\rP u^2\kappa\,\E\left[ \frac{1}{\Lambda^{-1}+\frac{u}{\tau}} \right].
\end{align}
% the quadratic form
%\begin{align}
%(u\sqrt{\kappa}\hba)^T\left(\Sigmab^{-1}+\frac{u}{\tau}\Iden\right)^{-1}(u\sqrt{\kappa}\hba) = u^2\kappa\,\frac{1}{p}\h^T\left(\Sigmab^{-1}+\frac{u}{\tau}\Iden\right)^{-1}\h
%\end{align}
%concentrates around its expectation \cite{?} $u^2\kappa\,\tr\left(\left(\Sigmab^{-1}+\frac{u}{\tau}\Iden\right)\right)^{-1}.$

\somm{Explain/justify PL(2)'ness of these functions. Clearly explain how A2 is used\dots}
In the last line, $\Lambda$ is a random variable as in Definition \ref{def:Xi}. \fx{Also, we used Assumption \ref{ass:mu} together with the facts that $\h$ is independent of $\Sigmab$ and  that the function $(x_1,x_2)\mapsto x_1^2(x_2^{-1}+u/\tau)^{-1}$ is $\rm{PL}(3)$ assuming $x_2$ is bounded (see Lemma \ref{lem:PLbdd} for proof).}
\ct{Question: Is it immediate that the empirical distribution of $(\beta,\Sigmab,\h)$ converges in $W_k$ to $\mu\otimes\Nn(0,1)$ given that $(\beta,\Sigmab)$ converges to $\mu$ and $\h$ is independent???}
%
 Second, we find that
\begin{align}
(\betas)^T\Sigmab^{-1/2}\left(\Sigmab^{-1}+\frac{u}{\tau}\Iden\right)^{-1}\Sigmab^{-1/2}\betas &= \frac{1}{p}(\sqrt{p}\betas)^T\left(\Iden+\frac{u}{\tau}\Sigmab\right)^{-1}(\sqrt{p}\betas) \nn\\
&= \frac{1}{p}\sum_{i=1}^{p}\frac{\left(\sqrt{p}\betas_i/\sqrt{\lai}\right)^2}{\lai^{-1}+\frac{u}{\tau}}\nn\\
&\rP \E\left[\frac{B^2\Lambda^{-1}}{\Lambda^{-1}+\frac{u}{\tau}}\right].
\end{align}
Here, $\Lambda,B$ are random variables as in Definition \ref{def:Xi} and \fx{we also used Assumption \ref{ass:mu} together with the fact that the function $(x_1,x_2)\mapsto x_1^2x_2^{-1}(x_2^{-1}+u/\tau)^{-1}$ is $\rm{PL}(2)$ assuming $x_2$ is bounded (see Lemma \ref{lem:PLbdd} for proof).}
Third, \cts{by independence of $(\betas, \Sigmab)$ from $\h$}
\begin{align}
(u\sqrt{\kappa}\hba)^T\left(\Sigmab^{-1}+\frac{u}{\tau}\Iden\right)^{-1}\Sigmab^{-1/2}\betas = u\sqrt{\kappa} \cdot \frac{1}{p}\sum_{i=1}
^p \frac{\h_i\Sigmab_{i,i}^{-1/2}(\sqrt{p}\betas_i)}{\Sigmab_{i,i}^{-1}+\frac{u}{\tau}\Iden} ~\rP~ 0.
\end{align}
Putting these together, the objective $\Rc(u,\tau)$ in \eqref{eq:AO_4} converges point-wise in $u,\tau$ to
\begin{align}
\Rc(u,\tau)\rP\Dc(u,\tau) := \frac{1}{2}\left({u\tau} + \frac{u\sigma^2}{\tau} - u^2\kappa\,\E\left[ \frac{1}{\Lambda^{-1}+\frac{u}{\tau}} \right] - \E\left[\frac{B^2\Lambda^{-1}}{\Lambda^{-1}+\frac{u}{\tau}}\right] \right).\label{eq:conv_pt}
\end{align}
Note that $\Rc(u,\tau)$ (and thus, $\Dc(u,\tau)$) is convex in $\tau$ and concave in $u$. Thus, the convergence in \eqref{eq:conv_pt} is in fact uniform (e.g., \cite{AG1982}) and we can conclude that 
\begin{align}\label{eq:Dc0}
\phi(\g,\h) \rP \max_{u\geq 0}\min_{\tau\in\Tc}~ \Dc(u,\tau).
\end{align}
and \fx{using strict concave/convexity of $\Dc(u,\tau)$, we also have the parameter convergence}
\somm{This line follows from strict convex/concavity right?}
\begin{align}\label{eq:Dc}
\fx{(u_n,\tau_n) \rP (u_*,\tau_*):=\arg\max_{u\geq 0}\min_{\tau\in\Tc}~ \Dc(u,\tau).}
\end{align}
In the proof of statement (iii) below, we show that the saddle point of \eqref{eq:Dc0} is $(u_*,\tau_*)$. In particular, $\tau_*$ is strictly in the interior of $\Tc$, which combined with convexity implies that
$$
 \max_{u\geq 0}\min_{\tau\in\Tc}~ \Dc(u,\tau) =  \max_{u\geq 0}\min_{\tau>0}~ \Dc(u,\tau) =: \bar\phi.
$$
This, together with the first display above proves the second statement of the lemma.

\vspace{5pt}
\subsubsection{Proof of (iii).}  Next, we compute the saddle point $(u_*,\tau_*)$ by studying the first-order optimality conditions of the strictly concave-convex $\Dc(u,\tau)$. Specifically, we consider the unconstrained minimization over $\tau$ and we will show that the minimum is achieved in the strict interior of $\Tc$. Direct differentiation of $\Dc(u,\tau)$ gives
\begin{subequations}
\begin{align}
{\tau} + \frac{\sigma^2}{\tau} - 2u\kappa\E\left[ \frac{1}{\Lambda^{-1}+\frac{u}{\tau}} \right] + \frac{u^2}{\tau}\kappa \E\left[ \frac{1}{\left(\Lambda^{-1}+\frac{u}{\tau}\right)^2}\right] + \frac{1}{\tau}\E\left[\frac{B^2\Lambda^{-1}}{\left(\Lambda^{-1}+\frac{u}{\tau}\right)^2}\right] &= 0, \label{eq:fo1}\\
{u} - \frac{u\sigma^2}{\tau^2} - \frac{u^3}{\tau^2}\kappa \E\left[ \frac{1}{\left(\Lambda^{-1} + \frac{u}{\tau}\right)^2}\right] - \frac{u}{\tau^2} \E\left[\frac{B^2\Lambda^{-1}}{\left(\Lambda^{-1}+\frac{u}{\tau}\right)^2}\right] &= 0,\label{eq:fo2}
\end{align}
\end{subequations}
Multiplying \eqref{eq:fo2}  with $\frac{\tau}{u}$ and adding to \eqref{eq:fo1} results in the following equation
\begin{align}
\tau = u\kappa\E\left[ \frac{1}{\Lambda^{-1}+\frac{u}{\tau}} \right] ~\Leftrightarrow ~
\E\left[ \frac{1}{(\frac{u}{\tau}\Lambda)^{-1}+1} \right] = \frac{1}{\kappa} \label{eq:tauu}\,.
\end{align}
Thus, we have found that the ratio $\frac{u_*}{\tau_*}$ is the unique solution to the equation in \eqref{eq:tauu}. Note that this coincides with the Equation \eqref{eq:ksi} that defines the parameter $\xi$ in Definition \ref{def:Xi}.  The fact that \eqref{eq:tauu} has a unique solution for all $\kappa>1$ can be easily seen as $F(x)=\E\left[ \frac{1}{(x\Lambda)^{-1}+1} \right], x\in\R_+$ has range $(0,1)$ and is strictly increasing (by differentiation). 


Thus, we call $\xi=\frac{u_*}{\tau_*}$. Moreover, multiplying \eqref{eq:fo2} with $u$ leads to the following equation for $\tau_*$:
\begin{align}
u_*^2 = \sigma^2\xi^2 + u_*^2 \xi^2  \kappa \E\left[ \frac{1}{\left(\Lambda^{-1} +\xi\right)^2}\right] + \xi^2  \E\left[ \frac{B^2\Lambda^{-1}}{\left(\Lambda^{-1} + \xi\right)^2}\right] ~\Rightarrow~ \tau_*^2 = \frac{\sigma^2 + \E\left[ \frac{B^2\Lambda^{-1}}{\left(\Lambda^{-1} + \xi\right)^2}\right]}{1-\xi^2\kappa \E\left[ \frac{1}{\left(\Lambda^{-1} +\xi\right)^2}\right]} =  \frac{\sigma^2 + \E\left[ \frac{B^2\Lambda^{-1}}{\left(\Lambda^{-1} + \xi\right)^2}\right]}{1-\kappa \E\left[ \frac{1}{\left((\xi\Lambda)^{-1} +1\right)^2}\right]}.
\end{align}
{Again, note that this coincides with Equation \eqref{eq:gamma} that determines the parameter $\gamma$ in Definition \ref{def:Xi}, i.e., $\tau_*^2 = \gamma.$ }

\vspace{5pt}
\subsubsection{Proof of (iv).}  For convenience, define 
$$F_n(\bth,\betas,\Sigmab):=\frac{1}{p} \sum_{i=1}^{p} f\left(\sqrt{p}\bth_{n,i},\sqrt{p}\betas_{i},\Sigmab_{ii}\right)\quad\text{and}\quad\alpha_*:=\E_\mu\left[f(X_{\kappa,\sigma^2}(\Lambda,B,H),B,\Lambda)\right].$$
Recall from \eqref{eq:w_n} the explicit expression for $\wh_n$, repeated here for convenience.
\begin{align}
\wh_n = -\left(\Sigmab^{-1}+\frac{u_n}{\tau_n}\Iden\right)^{-1}\left(\Sigmab^{-1/2}\betas-u_n\sqrt{\kappa}\hba\right).\nn
\end{align}
Also, recall that $\bth_n = \Sigmab^{-1/2}\wh_n+\betas$. Thus,  (and using the fact that $\hba$ is distributed as $-\hba$),
\begin{align}
\bth_n &=\Sigmab^{-1/2}\left(\Sigmab^{-1}+\frac{u_n}{\tau_n}\Iden\right)^{-1}u_n\sqrt{\kappa}\hba + \left(\Iden-\Sigmab^{-1/2}\left(\Sigmab^{-1}+\frac{u_n}{\tau_n}\Iden\right)^{-1}\Sigmab^{-1/2}\right)\betas\nn\\
\Longrightarrow\bth_{n,i}&=\frac{{\Sigmab_{i,i}^{-1/2}}}{1+(\ksi_n\Sigmab_{i,i})^{-1}}\sqrt{\kappa}\tau_n\hba_i +\left(1-\frac{1}{1+\ksi_n\Sigmab_{i,i}}\right)\betas_i.\label{eq:betan}
\end{align}
For $i\in[p]$, define 
\begin{align}
\vb_{n,i} = \frac{{\Sigmab_{i,i}^{-1/2}}}{1+(\ksi_*\Sigmab_{i,i})^{-1}}\sqrt{\kappa}\tau_*\hba_i +\left(1-\frac{1}{1+\ksi_*\Sigmab_{i,i}}\right)\betas_i
%-\left(\Sigmab^{-1}+\frac{u_*}{\tau_*}\Iden\right)^{-1}\left(\Sigmab^{-1/2}\betas-u_*\sqrt{\kappa}\hba\right).
\end{align}
In the above, for convenience, we have denoted $\xi_n:=u_n/\tau_n$ and recall that $\xi_*:=u_*/\tau_*$. 


The proof proceeds in two steps. In the first step, we use the fact that $\ksi_n\rP\ksi_*$ and $u_n\rP u_\star$ (see \eqref{eq:Dc}) to prove that for any $\eps\in(0,\ksi_*/2)$, there exists an absolute constant $C>0$ such that wpa 1:
\begin{align}\label{eq:ivstep1}
|F_n(\sqrt{p}\bth_n,\sqrt{p}\betas,\Sigmab) - F_n(\sqrt{p}\vb_n,\sqrt{p}\betas,\Sigmab) | \leq C\eps.
\end{align}
In the second step, we use \fx{pseudo-Lipschitzness of $f$ and Assumption \ref{ass:mu} to prove that}
\somm{Doesn't Assump 3 come out of nowhere? So far it looks like we need A1 and A2. We need to ensure consistency.} 
\begin{align}
|F_n(\vb_n,\betas,\Sigmab)| \rP \alpha_*.\label{eq:ivstep2}
\end{align}
The desired follows by combining \eqref{eq:ivstep1} and \eqref{eq:ivstep2}. Thus, in what follows, we prove \eqref{eq:ivstep1} and \eqref{eq:ivstep2}.


\vspace{2pt}
\noindent\underline{Proof of \eqref{eq:ivstep1}.}~~Fix some $\eps\in(0,\ksi_*/2)$. From \eqref{eq:Dc}, we know that w.p.a. 1 $|\xi_n-\xi_*|\leq \eps$ and $|u_n-u_*|\leq \eps$. Thus, $\wh_n$ is close to $\vb_n$. Specifically, in this event, for every $i\in[p]$, it holds that:
\somm{Fourth line: Constant missing here. Powers are wrongish}
\begin{align}
\nn|\bth_{n,i} - \vb_{n,i}| &\leq {|\betas_i|}\left|\frac{1}{1+\xi_n\lai}-\frac{1}{1+\xi_*\lai}\right| + \sqrt{\kappa}\frac{|\hba_i|}{\sqrt{\Sigmab_{i,i}}}\left|\frac{\tau_n}{1+(\xi_n\lai)^{-1}}-\frac{\tau_*}{1+(\xi_*\lai)^{-1}}\right| \\
\nn&= {|\betas_i|}\left|\frac{1}{1+\xi_n\lai}-\frac{1}{1+\xi_*\lai}\right| + \sqrt{\kappa}\frac{|\hba_i|}{\sqrt{\Sigmab_{i,i}}}\left|\frac{u_n}{\xi_n+\lai^{-1}}-\frac{u_*}{\xi_*+\lai^{-1}}\right| \\
\nn& \leq
{|\betas_i|}\frac{|\lai||\xi_n-\xi_*|}{|1+\xi_n\lai||1+\xi_*\lai|} + \sqrt{\kappa}\frac{|\hba_i|}{\sqrt{\lai}}  \frac{u_*|\ksi_n-\ksi_*|}{(\ksi_n+\lai^{-1})(\ksi_*+\lai^{-1})} + 
\sqrt{\kappa}\frac{|\hba_i|}{\sqrt{\lai}}\frac{|u_n-u_*|}{\xi_n+\lai^{-1}}
\\
\nn& \leq
{|\betas_i|}{\Sigma_{\max}\eps} + \sqrt{\kappa}{|\hba_i|} u_*\Sigma_{\max}^{3/2} \eps+ 
\sqrt{\kappa}|\hba_i|\Sigma_{\max}^{1/2}\eps
\\
&\leq \eps \cdot \max\left\{\Sigma_{\max}^{3/2},\Sigma_{\max}^{1/2}\right\}\, \left(|\hba_i| + |\betas_i|\right).\label{eq:betav}
%\nn& \leq
%\frac{|\betas_i|}{\sqrt{\lai}} \frac{\eps}{(\xi-\eps)\xi} + \sqrt{\kappa}|\hba_i|\left(\frac{\eps(u_*+\eps)}{(\xi-\eps)\xi} + \frac{\eps}{\xi}\right)\\
%& \leq
%\frac{|\betas_i|}{\sqrt{\lai}} c_1(\eps) + |\hba_i|c_2(\eps).\label{eq:wivi}
\end{align}%\som{$\eps$ missing}
In the second line above, we recalled that $u_n=\tau_n\xi_n$ and $u_*=\tau_*\xi_*$. In the third line, we used the triangle inequality. In the fourth line, we used that $\ksi_*>0$, $0<\lai\leq\Sigma_{\max}$ and $\ksi_n\geq \ksi_*-\eps \geq \ksi_*/2 >0$. 
%In the second line above, we used the triangle inequality. In the third inequality we used Assumption \ref{ass:inv} that $\lai>0, i\in[p]$, and in the last line we defined appropriate constants $c_1(\eps),c_2(\eps)>0$. Importantly,  ... 

Now, we will use this and Lipschitzness of $f$ to argue that there exists absolute constant $C>0$ such that wpa 1:
\begin{align}
|F_n(\sqrt{p}\bth_n,\sqrt{p}\betas,\Sigmab) - F_n(\sqrt{p}\vb_n,\sqrt{p}\betas,\Sigmab) | \leq C\eps.\nn
\end{align}

Denote, $\ab_i=(\sqrt{p}\bth_{n,i},\sqrt{p}\betas_i,\Sigmab_{i,i})$ and $\bb_i=(\sqrt{p}\vb_{n,i},\sqrt{p}\betas_i,\Sigmab_{i,i})$. Following the exact same argument as in \eqref{eq:dev2show} (just substitute $\bt\leftrightarrow\vb_n$ in the derivation), we have that for some absolute constant $C>0$ wpa 1:
\begin{align}
|F_n(\bth_n(\g,\h),\betas,\Sigmab) - F_n(\vb_n,\betas,\Sigmab)| &\leq C \|\bth_n-\vb_n\|_2.
\end{align}
%\begin{align}
%|F_n(\bth_n(\g,\h),\betas,\Sigmab) - F_n(\vb_n,\betas,\Sigmab)| &\leq \frac{L}{p}\sum_{i=1}^p (1+\max\{\|\ab_i\|_2^{k-1},\|\bb_i\|_2^{k-1}\})\|\ab_i-\bb_i\|_2\nn \\
%&\leq \frac{L}{p}\sum_{i=1}^p \left(1+\max\{\|\ab_i\|_2^{k-1},\|\bb_i\|_2^{k-1}\}\right)\cdot|\sqrt{p}\bth_{n,i}(\g,\h) - \sqrt{p}\vb_{n,i}|\nn\\
%&\leq L\left(1+\max\{\frac{1}{p}\sum_{i=1}^p \|\ab_i\|_2^{2k-2},\frac{1}{p}\sum_{i=1}^p\|\bb_i\|_2^{2k-2}\} \right)^{1/2}{\|\bth_n(\g,\h) - \vb_n\|_2}\nn.
%\end{align}
From this and \eqref{eq:betav}, we find that
\begin{align}
|F_n(\bth_n(\g,\h),\betas,\Sigmab) - F_n(\vb_n,\betas,\Sigmab)|
&\leq  
C\eps\max\left\{\Sigma_{\max}^{3/2},\Sigma_{\max}^{1/2} \right\} \left(\sum_{i=1}^p \left(|\hba_i|+|\betas_i|\right)^2\right)^{1/2}\nn
\\
 &\leq 
C\eps\sqrt{2}\max\left\{\Sigma_{\max}^{3/2},\Sigma_{\max}^{1/2}\right\}\sqrt{\|\betas\|_2^2 + \|\hba\|_2^2}.\label{eq:epsS}
\end{align}
%As in \eqref{eq:bdmom}, it can be shown that wpa 1: $\frac{1}{p}\sum_{i=1}^p \|\ab_i\|_2^{2k-2}<\infty$ and $\frac{1}{p}\sum_{i=1}^p \|\bb_i\|_2^{2k-2}<\infty$, as $p\rightarrow\infty$. 
But, \cts{recall that $\|\betas\|_2^2 = \frac{1}{p}\sum_{i=1}^p(\sqrt{p}\betas_i)^2<\infty$, as $p\rightarrow \infty$ by Assumption \ref{ass:mu}.}
%{\color{red} by assumption on second moment convergence of $\sqrt{p}\betas$.} 
Also, since $\hba_i\sim\Nn(0,1/p)$, it holds that $\|\hba\|_2^2\leq 2$, wpa 1 as $p\rightarrow\infty$. Therefore, from \eqref{eq:epsS}, wpa 1, there exists constant $C>0$ such that 
\begin{align}
|F_n(\bth_n(\g,\h),\betas,\Sigmab) - F_n(\vb_n,\betas,\Sigmab)| \leq C \cdot\eps,
\end{align}
as desired.


%
%
%
% By $f\in\rm{PL}(k)$, we have that
%\begin{align}
%|f(\sqrt{p}\bth_{n,i},\sqrt{p}\betas_i,\Sigmab_{i,i}) - f(\sqrt{p}\vb_{n,i},\sqrt{p}\betas_i,\Sigmab_{i,i}) | \leq L ( 1 +  \|\ab_i\|_2^{k-1} +  \|\bb_i\|_2^{k-1} ) \sqrt{p} |\bth_{n,i}-\vb_{n,i}|. \nn
%\end{align}•
%Thus, using Cauchy-Swchartz,
%\begin{align}
%|F_n(\sqrt{p}\bth_n,\sqrt{p}\betas,\Sigmab) - F_n(\sqrt{p}\vb_n,\sqrt{p}\betas,\Sigmab) | &\leq \frac{L}{p} \sum_{i=1}^p ( 1 +  \|\ab_i\|_2^{k-1} +  \|\bb_i\|_2^{k-1} ) \sqrt{p} |\bth_{n,i}-\vb_{n,i}|. \\
%\end{align}•
%
% Using $f\in\rm{PL}(2)$, this implies the following
%\begin{align}
%|\frac{1}{p}\sum_{i=1}^{p}f(\sqrt{p}\w_{n,i})-\frac{1}{p}\sum_{i=1}^{p}f(\sqrt{p}\vb_{n,i})| &\leq \frac{1}{p}\sum_{i=1}^{p} L(1+\sqrt{p}|\w_{n,i}|+\sqrt{p}|\vb_{n,i}|)\sqrt{p}|\w_{n,i}-\vb_{n,i}| \nn
%\\
%&\leq \nn
%\frac{L}{p}\sum_{i=1}^{p}(1+\sqrt{p}|\w_{n,i}|+\sqrt{p}|\vb_{n,i}|) \left(|\nub_i| c_1 + |\h_i|c_2\right)\nn\\
%&\leq C\left(\frac{\tn{\nub}}{\sqrt{p}}+ \frac{\tn{\h}}{\sqrt{p}}\right)(1+B+\tn{\vb_{n}}) \label{eq:fwivi}
%%&\leq C(1+2B)\left(\frac{\tn{\nub}}{\sqrt{p}}+ \frac{\tn{\h}}{\sqrt{p}}\right) 
%\end{align}\som{$\eps$ missing}
%The inequality in the first line follows from $f\in\rm{PL}(2)$. The inequality in the second line uses \eqref{eq:wivi}, and we have also recalled the notation $\nub_i=\sqrt{p}\betas_i/\sqrt{\lai}$ in Assumption \ref{ass:mu} and $\h_i=\sqrt{p}\hba_i$. The inequality in the third line uses Cauchy-Swchartz, $\tn{\w_{n}}\leq B$ and $C:=C(L,c_1,c_2)$ is a universal constant.
%
%By Assumption \ref{ass:mu}, $\tn{\nub}/{\sqrt{p}}\rP\sqrt{\E[N^2]}<\infty$. Also, with probability approaching 1: $\h/{\sqrt{p}}<2$. Similarly, it can be argued that with probability approaching 1 $\tn{\wh_n}$ is bounded by a large enough constant. Combined with \eqref{eq:fwivi}, we have shown that there exists universal constant $C'(\eps)>0$ such that with probability approaching 1:
%\begin{align}
%|\frac{1}{p}\sum_{i=1}^{p}f(\sqrt{p}\w_{n,i})-\frac{1}{p}\sum_{i=1}^{p}f(\sqrt{p}\vb_{n,i})| \leq C'(\eps) .
%\end{align}


\vspace{2pt}
\noindent\underline{Proof of \eqref{eq:ivstep2}.}~~Next, we will use \ref{ass:mu} to show that 
\begin{align}
|F_n(\vb_n,\betas,\Sigmab)| \rP \alpha_*.\label{eq:AO_conv}
\end{align}
Notice that $\vb_n$ is a function of $\betas,\Sigmab,\hba$. Concretely, define ${\ggt}:\R^3\rightarrow\R$, such that
$$
\ggt(x_1,x_2,x_3) := \frac{x_2^{-1/2}}{1+(\ksi_* x_2)^{-1}}\sqrt{\kappa}\tau_* x_3 + (1-(1+\ksi_*x_2)^{-1})x_1,
$$
and notice that
$$
\sqrt{p}\vb_{n,i} = \ggt\left(\sqrt{p}\betas_i,\lai,\h_i\right) = \frac{{\Sigmab_{i,i}^{-1/2}}}{1+(\ksi_*\Sigmab_{i,i})^{-1}}\sqrt{\kappa}\tau_*\h_i +\left(1-\frac{1}{1+\ksi_*\Sigmab_{i,i}}\right)\sqrt{p}\betas_i.
$$
Thus,
$$
F_n(\vb_n,\betas,\Sigmab) = \frac{1}{p}\sum_{i=1}^p f\left(g\left(\sqrt{p}\betas_i,\lai,\h_i\right),\sqrt{p}\betas_i,\lai\right) =: \frac{1}{p}\sum_{i=1}^p h\left(\h_i,\sqrt{p}\betas_i,\lai\right),
$$
where we have defined $h:\R^3\rightarrow\R$:
\begin{align}
h(x_1,x_2,x_3) := f\left(\ggt(x_2,x_3,x_1),x_2,x_3\right).\label{h func}
\end{align}
\cts{We will prove that $h\in\rm{PL}(4)$. Indeed, if that were the case, then Assumption \ref{ass:mu} gives}
\begin{align}
\frac{1}{p}\sum_{i=1}^p h\left(\h_i,\sqrt{p}\betas_i,\lai\right) \rP \E_{\Nn(0,1)\otimes \mu}\left[h(H,B,\Lambda)\right] &=  \E_{\Nn(0,1)\otimes \mu}\left[f\left(\ggt(B,\Lambda,H),B,\Lambda\right))\right]\\
& = \E[f\left(X_{\kappa,\sigma^2}(\Lambda,B,H),B,\Lambda\right)] = \alpha_*,
\end{align}
where the penultimate equality follows by recognizing that (cf. Eqn. \eqref{eq:X})
$$
\ggt(B,\Lambda,H) = (1-(1+\ksi_*\Lambda)^{-1})B + \sqrt{\kappa}\frac{\tau_*\Lambda^{-1/2}}{1+(\ksi_*\Lambda)^{-1}}H =  X_{\kappa,\sigma^2}(\Lambda,B,H).
$$
It remains to show that $h\in\rm{PL}(4)$. Lemma \ref{lem:hPL} in Section \ref{SM useful fact} shows that if $f\in\rm{PL}(k)$, then $h\in\rm{PL}(k+1)$ for all integers $k\geq 2$.  First, consider the case $f\in\rm{PL}(2)$. Then,  $h\in\rm{PL}(3)$; thus, also $h\in\rm{PL}(4)$. Second, consider the case $f=f_\Lc$. Then, we prove in Lemma \ref{lem:fL}  that $f_\Lc\in\rm{PL}(3)$. Thus, the desired holds in this case too.  
%In the remaining of the proof, \cts{we show that $h\in\rm{PL}(k)$} which is formalized below.

\vspace{5pt}
\subsubsection{Proof of (v):} Let $\psi:\R\rightarrow\R$ be any bounded Lipschitz function. The function $f(a,b,c) = \psi(a)$ is trivially $\rm{PL}$ of order $2$. Thus, by directly applying statement (iv) of the lemma, we find that
$$
\frac{1}{p}\sum_{i=1}^p{\psi(\sqrt{p}\bth_{n,i}(\g,\h))} \rP \E\left[\psi(X_{\kappa,\sigma^2})\right].
$$
Since this holds for any bounded Lipschitz function, we have shown that the empirical convergence of $\bth_n$ converges weakly to the distribution of $X_{\kappa,\sigma^2}$. It remains to prove boundedness of the $2$nd moment as advertised in \eqref{eq:k_AO}. Recall  from
\eqref{eq:betan} that
\begin{align}
\sqrt{p}\bth_{n,i}&=\frac{{\Sigmab_{i,i}^{-1/2}}}{1+(\ksi_n\Sigmab_{i,i})^{-1}}\sqrt{\kappa}\tau_n\h_i +\left(1-\frac{1}{1+\ksi_n\Sigmab_{i,i}}\right)(\sqrt{p}\betas_i).\nn
\end{align}
Using this, boundedness of $\Sigmab_{i,i}$ from Assumption \ref{ass:mu}, and the fact that $\tau_n\rP\tau_\star, \ksi_n\rP\ksi_\star$, there exists constant $C=C(\Sigma_{\max},\Sigma_{\min},k,\tau_\star,\ksi_\star)$ such that wpa 1,
\begin{align}
\frac{1}{p}\sum_{i=1}^p|\sqrt{p}\bth_{n,i}|^{2} \leq C\left(\frac{1}{p}\sum_{i=1}^p|\h_{i}|^{2}+\frac{1}{p}\sum_{i=1}^p|\sqrt{p}\betas_{i}|^{2}\right).\nn
\end{align}
But the two summands in the expression above are finite in the limit of $p\rightarrow\infty$. Specifically, (i) from Assumption \ref{ass:mu}, $\frac{1}{p}\sum_{i=1}^p|\sqrt{p}\betas_{i}|^{2}\rP\E[B^{2}]<\infty$; (ii) $\frac{1}{p}\sum_{i=1}^p|\h_{i}|^{2}\rP\E[H^{2}]=1$, using the facts that $\h_i\simiid\Nn(0,1)$ and $H\sim\Nn(0,1)$. This proves \eqref{eq:k_AO}, as desired.

%%
%%
%%\begin{lemma} The function $h$ in \eqref{h func} is $\rm{PL}(k)$ as long as $f$ is $\rm{PL}(k)$.
%%{\color{red}The function $h$ is $\R^3\rightarrow\R$, not $\R^{3p}\rightarrow\R$! See Lemma 6 below.}
%%\end{lemma}
%%\begin{proof} 
%%Since $f$ is $\rm{PL}(k)$, for some $L>0$, \eqref{PL func} holds. Fix $\ab=[\x,\y,\z]\in\R^{3p},\ab'=[\x',\y',\z']\in\R^{3p}$. Let $\bb=[\g,\y,\z]$ where $\g=g(\y,\z,\x)$. We have that
%%\begin{align}
%%|h(\ab)-h(\ab')|&\leq |f(\bb)-f(\bb')|\\
%%&\leq L\left(1+\|\bb\|_2^{k-1}+\|\bb'\|_2^{k-1}\right)\|\bb-\bb'\|_2\\
%%&\lesssim L\left(1+\|\ab\|_2^{k-1}+\|\ab'\|_2^{k-1}+\|\g\|_2^{k-1}+\|\g'\|_2^{k-1}\right)(\|\ab-\ab'\|_2+\|\g-\g'\|_2).
%%%~~~~~~~~&L\left(1+\|\g(\y,\z,\x)\|_2^{k-1}+\|g(\y',\z',\x')\|_2^{k-1}\right)\|g(\y,\z,\x)-g(\y',\z',\x')\|_2.
%%\end{align}
%%Next, we need to bound the $\g$ term in terms of $\ab$. This is accomplished as follows
%%\begin{align}
%%\|\g\|_2^{k-1}&\lesssim \frac{1}{p}\sum_{i=1}^p|g_i|^{k-1}\\
%%&\lesssim \frac{1}{p}\sum_{i=1}^p\left|\frac{y_i^{-1/2}}{1+(\ksi y_i)^{-1}}\sqrt{\kappa}\tau_* z_i + (1-(1+\ksi_*y_i)^{-1})x_i\right|\\
%%&\lesssim \frac{1}{p}\sum_{i=1}^p(|z_i|+|x_i|)^{k-1}\lesssim \|\x\|_2^{k-1}+\|\z\|_2^{k-1}\\
%%&\lesssim \|\ab\|^{k-1}.
%%\end{align}
%%Secondly and similarly, we have the following perturbation bound on the $\g-\g'$. \so{Suppose the triples $(x_i,y_i,z_i)$ takes values in a fixed bounded compact set $\Mc$. We will prove the following sequence of inequalities} 
%%\begin{align}
%%%|\|\g\|_2-\|\g'\|_2|&\leq 
%%\tn{\g-\g'}^2&\leq \sum_{i=1}^p|g(x_i,y_i,z_i)-g(x'_i,y'_i,z'_i)|^2\\
%%&\leq C^2_{\Mc} \sum_{i=1}^p(|x_i-x'_i|^2+|y_i-y'_i|^2+|z_i-z'_i|^2)\label{sec line eq}\\
%%&\leq C^2_{\Mc} (\tn{\x-\x'}^2+\tn{\y-\y'}^2+\tn{\z-\z'}^2)\\
%%&\lesssim \tn{\ab-\ab'}^2.
%%\end{align}
%%In what follows, we prove the second line \eqref{sec line eq} i.e.~the fact that for any triples $a=(x,y,z),a'=(x',y',z')$ (with $a,a'\in\R^3$), we have that 
%%\[
%%|g(a)-g(a')|\leq C_{\Mc}\tn{a-a'}.
%%\]
%%Set $C_\nabla=\sup_{a\in \Mc}\tn{\nabla g(a)}$. By definition of gradient, the inequality above holds with $C_\Mc=C_\nabla$. Thus, all that remains is proving that $C_\nabla$ is upper bounded by a constant. \so{However, this automatically holds because from the definition of $g(\dots)$ function, it is clear that $C_\nabla$ is defined everywhere and continuous thus it has a finite maximum over a compact set.} \cts{Only the problem is that $z_i$ (i.e., $\h_i$) is $\Nn(0,1)$ so it is not bounded.}
%%% is bounded by a constant.
%%\end{proof}
%%


%But from Assumption \ref{ass:mu} and $f\in\rm{PL(2)}$ it follows immediately that with probability approaching 1,
%\begin{align}
%|\frac{1}{p}\sum_{i=1}^{p}f(\sqrt{p}\vb_{n,i})-\E_\mu\left[X_{\kappa,\sigma^2}(\Lambda,N,H)  \right]| \leq \eps.
%\end{align}
%Combining the above displays and using triangle inequality completes the proof of the statement.




\section{Asymptotic formulas on Magnitude- and Hessian- pruning}\label{sec prune risk}
Here, we use Theorem \ref{thm:master_W2} to characterize the risk of the magnitude- and Hessian- pruned solutions. This section supplements the discussion in Section \ref{sec:risk}. For completeness, first, we recall the magnitude-pruning results of Section \ref{sec:risk} and restate Corollary \ref{cor:mag} below. This corollary characterizes the performance of magnitude pruning. Following this, we shall further discuss Hessian pruning.


\subsection{Magnitude-based pruning}

We begin with the following necessary definitions.  Define the hard-thresholding function with fixed threshold $t\in\R_+$ as follows:
\begin{align}\label{eq:threshold_app}
\Tc_t(x) = \begin{cases} x & \text{if } |x|>t \\ 0 & \text{otherwise} \end{cases}.
\end{align}
Further, {given model sparsity target $1>\alpha>0$}, define the threshold $t^\star$ as follows:
\begin{align}\label{eq:tstar_app}
t^\star:=\sup\left\{t\in\R\,:\, \Pr(|X_{\kappa,\sigma^2}|\geq t) \geq \alpha \right\}.
\end{align}


\cormag*

The proof of the corollary above, is given in Section \ref{sec:risk}. Below, we extend the results to Hessian-based pruning.


\subsection{Hessian-based pruning}
Let $\bth$ be the min-norm solution in \eqref{eq:min_norm}. Recall that the Hessian-pruned model (via Optimal Brain Damage) $\betab_s^H$ at sparsity $s$ is given by
\begin{align}\label{eq:hp}
\betab_s^H = \hat\Sigmab^{-1/2}\Tb_s(\hat\Sigmab^{1/2}\betab),
\end{align}
where $\hat\Sigmab=\diag{\X^T\X}/n$ the diagonal of the empirical covariance matrix. 

We will argue that the following formula characterizes the asymptotic risk of the Hessian pruning solution. 
Recall the notation in \eqref{eq:threshold_app} and define
\begin{align}\label{eq:tdiam_app}
t^\diamond:=\sup\left\{t\in\R\,:\, \Pr(|\Lambda^{1/2}X_{\kappa,\sigma^2}|\geq t) \geq \alpha \right\}.
\end{align}
The risk of the Hessian-pruned model satisfies the following in the limit of $n,p,s\rightarrow\infty$ at rates $\kappa:=p/n>1$ and $\alpha:=s/p\in(0,1)$ (cf. Assumption \ref{ass:linear}):
\begin{align}\label{eq:hessian_risk}
\Lc(\bth^H_s) \rP \sigma^2 + \E\left[ \left(\Lambda^{1/2}B-\Tc_{t^\diamond}(\Lambda^{1/2}X_{\kappa,\sigma^2})\right)^2 \right],
\end{align}
where the expectation is over $(\Lambda,B,H)\sim\mu\otimes\Nn(0,1)$. In our LGP experiments, we used this formula \eqref{eq:hessian_risk} to accurately predict the Hessian-based pruning performance. 
%\begin{corollary}[Risk of Hessian-pruning]\label{cor:hes}
%Let the same assumptions and notation as in the statement of Theorem \ref{thm:master_W2} hold. Specifically, let $\hat\betab$ be the min-norm solution in \eqref{eq:min_norm} and $\hat\beta_s^H$ the Hessian-pruned model at sparsity $s$. Recall the notation in \eqref{eq:threshold_app} and define
%\begin{align}\label{eq:tdiam_app}
%t^\diamond:=\inf\left\{t\in\R\,:\, \Pr(|\Lambda^{1/2}X_{\kappa,\sigma^2}|\geq t) \geq \alpha \right\}.
%\end{align}
%The risk of the magnitude-pruned model satisfies the following in the limit of $n,p,s\rightarrow\infty$ at rates $\kappa:=p/n>1$ and $\alpha:=s/p\in(0,1)$ (cf. Assumption \ref{ass:linear}):
%\begin{align}
%\Lc(\bth^H_s) \rP \sigma^2 + \E\left[ \Lambda\left(B-\Lambda^{-1/2}\Tc_{t^\diamond}(\Lambda^{1/2}X_{\kappa,\sigma^2})\right) \right],\nn
%\end{align}
%where the expectation is over $(\Lambda,B,H)\sim\mu\otimes\Nn(0,1).$
%\end{corollary}
%\begin{proof}

Recall the definition of the hard-thresholding operator $\Tc_t(x)$. Similar to Section \ref{sec:risk}, we consider a threshold-based pruning vector $$\bth^{\Tc,H}_{t} := \hat\Sigmab^{-1/2}\Tc_{t/\sqrt{p}}(\hat\Sigmab^{1/2}\hat\bt),$$ where $\Tc_t$ acts component-wise. Further define 
$$\bth^{\Tc,H^\star}_{t} := \Sigmab^{-1/2}\Tc_{t/\sqrt{p}}(\Sigmab^{1/2}\hat\bt).$$
Note that $\bth^{\Tc,H^\star}_{t}$ uses the true (diagonal) covariance matrix $\Sigmab$ instead of its sample estimate $\hat\Sigmab$. For later reference, note here that $\hat\Sigmab$ concentrates (entry-wise) to $\Sigmab$. Using boundedness of $\Sigmab$ and standard concentration of sub-exponential random variables.


First, we compute the limiting risk of $\bth^{\Tc,H^\star}_{t}$. Then, we will use the fact that $\hat\Sigmab$ concentrates (entry-wise) to $\Sigmab$, to show that the risks of $\bth^{\Tc,H^\star}_{t}$ and $\bth^{\Tc,H}_{t}$ are arbitrarily close as $p\rightarrow\infty$.

%We start by computing the risk of $\bth^{\Tc,H^\star}_{t}$. 
Similar to \eqref{eq:loss_w2}, 
\begin{align}
\Lc(\bth^{\Tc,H^\star}_t) &= \sigma^2 + (\bt^\star-\bth^{\Tc,H^\star}_t)^T\Sigmab(\bt^\star-\bth^{\Tc,H^\star}_t) \nn \\
  &= \sigma^2 + \frac{1}{p}\sum_{i=1}^p\Sigmab_{i,i}\big(\sqrt{p}\betas_i-\lai^{-1/2}\Tc_{t}(\lai^{1/2}\sqrt{p}\hat\betab_i)\big)^2 \nn
    = \sigma^2 + \frac{1}{p}\sum_{i=1}^p\big(\lai^{1/2}\sqrt{p}\betas_i-\Tc_{t}(\lai^{1/2}\sqrt{p}\hat\betab_i)\big)^2 
      \\
    &= \sigma^2 + \frac{1}{p}\sum_{i=1}^p\big(\lai^{1/2}\sqrt{p}\betas_i-\Tc_{t}(\wh_i + \lai^{1/2}\sqrt{p}\betas_i\big)^2 \nn \\
        &=: \sigma^2 + \frac{1}{p}\sum_{i=1}^p\big(\sqrt{p}\blat_i-\Tc_{t}(\wh_i + \blat_i)\big)^2.\label{eq:hesa}
  %\nn\\
%  &\rP \sigma^2 + \E\left[ \Lambda\left(B-\Tc_t(X_{\kappa,\sigma^2})\right) \ri
\end{align}
In the second line above, we used that $\sqrt{p}\Tc_{t/\sqrt{p}}(x)=\Tc_{t}(\sqrt{p}x)$. In the third line, we recalled the change of variable in \eqref{eq:w}, i.e., $\wh$ is the solution to \eqref{eq:PO}. Finally, in the last line we defined $\blat:=\sqrt{\Sigmab}\betas$ (note that this is related to the saliency score $\bla$ defined in \eqref{saliency eq}).

To evaluate the limit of the empirical average in \eqref{eq:hesa}, we proceed as follows. First, we claim that the empirical distribution of $\sqrt{p}\blat$ converges weakly to the distribution of the random variable $B\sqrt{\Lambda}$, where $(\Lambda,B)\sim\mu$. \so{Note that this convergence is already implied by the proof of Theorem \ref{thm:master_W2} by setting the $g$ function to be zero in \eqref{eq:fdef}.}
%\somm{This part is repetitive, can be simplified.}
For an explicit proof, take any bounded Lipschitz test function $\psi:\R\rightarrow\R$ and call $\psi'(x,y):=\psi({\sqrt{x}}y)$. Then, 
\begin{align}
|\psi'(\lai,\sqrt{p}\betas_i)-\psi'(\lai',\sqrt{p}{\betas_i}')| &= |\psi(\sqrt{\lai}\sqrt{p}\betas_i)-\psi(\sqrt{\lai'}\sqrt{p}{\betas_i}')| \leq C |\sqrt{\lai}\sqrt{p}\betas_i - \sqrt{\lai'}\sqrt{p}{\betas_i}'| \nn
\\
&\leq
C |\sqrt{p}\betas_i - \sqrt{p}{\betas_i}'| + C'|\sqrt{p}{\betas_i}'| |\sqrt{\lai}- \sqrt{\lai'}|\nn  \\
&\leq
C (|\sqrt{p}\betas_i - \sqrt{p}{\betas_i}'| + |\sqrt{p}{\betas_i}'| |{\lai}- {\lai'}|)\nn  \\
&\leq
C (1 + \|[\sqrt{p}\betas_i,\lai]\|_2 +\|[\sqrt{p}{\betas_i}',\lai']\|_2 ) \sqrt{|\sqrt{p}\betas_i - \sqrt{p}{\betas_i}'|^2   + |{\lai}- {\lai'}|^2}\nn.
\end{align}
Thus, $\psi'$ is $\rm{PL(2)}$. Hence, from Assumption \ref{ass:mu},
\begin{align}
\frac{1}{p}\sum_{i=1}^p \psi(\sqrt{p}\blat_i) = \frac{1}{p}\sum_{i=1}^p \psi'(\lai,\sqrt{p}\betas_i) \rP  \E[\psi'(B,\sqrt{\Lambda})] =  \E[\psi(B\sqrt{\Lambda})]\label{eq:convvvv}
\end{align}
Besides, from Theorem \ref{thm:master_W2} applied for $f_\Lc(x,y,z)=zy^2$ (i.e., set $g$ the zero function in \eqref{eq:fdef}), we have that
$$
\frac{1}{p}\sum_{i=1}^p \sqrt{p}(\blat_i)^2 \rP  \E[B^2\Lambda].
$$
Therefore, convergence in \eqref{eq:convvvv} actually holds for any $\psi\in\rm{PL}(2)$. Next, observe that, $\wo$ of \eqref{eq:w_n} can be written in terms of $\blat$ via \begin{align}\label{eq:w_n2}
\w_n=\w'(\tau_n,u_n)& = -\left(\Iden+\frac{u_n}{\tau_n}\Sigmab\right)^{-1}\left(\Sigmab^{1/2}\betas-u_n\sqrt{\kappa}\Sigmab\hba\right)\\
&=-\left(\Iden+\frac{u_n}{\tau_n}\Sigmab\right)^{-1}\left(\blat-u_n\sqrt{\kappa}\Sigmab\hba\right).
\end{align}
After this observation, the convergence proof can be finalized by using a modified version of Assumption \ref{ass:mu} as follows.%in two ways. First approach is

{
\begin{assumption}[Empirical distribution for saliency] \label{ass:lambda} Set $\blat_i=\sqrt{\lai}\bts_i$. The joint empirical distribution of $\{(\lai,\blat_i)\}_{i\in[p]}$ converges in {Wasserstein-k} distance to a probability distribution $\mu=\mu(\Lambda,S)$ on $\R_{>0}\times\R$ for some \fx{$k\geq 4$}. That is
$
\frac{1}{p}\sum_{i\in[p]}\delta_{(\lai,\sqrt{p}\blat_i)} \stackrel{W_k}{\Longrightarrow} \mu.
$
%weakly to a probability distribution $\mu=\mu(\Lambda, S)$ on $\R_{>0}\times\R$ with bounded $4$th moment and assume that as $p\rightarrow\infty$, the $4$th moment of the empirical distribution converges to the $4$th moment of $\mu$, i.e.,
%$$
%\frac{1}{p}\sum_{i=1}^p\left(\sqrt{(\sqrt{p}\blat_i)^{2}+\lai^2}\right)^4 \rP \E_{(\Lambda,S)\sim\mu}\left[\left(\sqrt{\Lambda^2+S^2}\right)^4\right]<\infty.
%$$
\end{assumption}
}

Under this assumption, it can be verified that, the exact same proof strategy we used for magnitude-based pruning would apply to Hessian-based pruning by replacing $\bts$ with $\blat$. The reason is that the Hessian pruning risk is a $\text{PL}(4)$ function of $\h,\blat,\bSi$ (e.g.~can be shown in a similar fashion to Lemma \ref{lem:hPL}). Observe that Assumption \ref{ass:lambda} is a reasonable assumption given the naturalness of the saliency score. If we only wish to use the earlier Assumption \ref{ass:mu} rather than Assumption \ref{ass:lambda}, one can obtain the equivalent result by modifying \ref{ass:mu} to enforce a slightly higher order bounded moment and convergence condition. Finally, one needs to address the perturbation due to the finite sample estimation of the covariance. Note that, even if the empirical covariance doesn't converge to the population, its diagonal weakly converges to the population (as we assumed the population is diagonal). The (asymptotically vanishing) deviation due to the finite sample affects can be addressed in an identical fashion to the deviation analysis of $\tau_n,u_n$ at \eqref{eq:ivstep1} and \eqref{eq:betav}. While these arguments are reasonably straightforward and our Hessian pruning formula accurately predicts the empirical performance, the fully formal proof of the Hessian-based pruning is rather lengthy to write and does not provide additional insights. 
%Thus the associated formal theoretical statement is deferred to a future version.

\cmt{
if one assumes Assumption \ref{ass:mu} with the 
The overall function $f(\blat_i,$ becomes
\[
h(x,y,z)=f(g(y,z,x),y,z)=y(x-g(z))^2
\]}
%\cts{Using the assumption of the corollary that $\beta$}
%\end{proof}








%%%%%%%%%%%%%%%%%%%%%%%%%%%%%%%%%%%%
%%
%%
%%%%%%%%%%%%%%%%%%%%%%%%%%%%%%%%%%%%
\section{Useful results about pseudo-Lipschitz functions}\label{SM useful fact}
For $k\geq 1$ we say a function $f:\R^m\rightarrow\R$ is pseudo-Lipschitz of order $k$ and denote it by $f\in \rm{PL}(k)$ if there exists a cosntant $L>0$ such that, for all $\x,\y\in\R^m$:
\begin{align}
|f(\x)-f(\y)|\leq L\left(1+\|\x\|_2^{k-1}+\|\y\|_2^{k-1}\right)\|\x-\y\|_2.\label{PL func}
\end{align}
In particular, when $f\in\rm{PL}(k)$, the following properties hold:
\begin{enumerate}
\item There exists a constant $L'$ such that for all $\x\in\R^n$: $|f(\x)|\leq L'(1+\|\x\|_2^k).$

\item $f$ is locally Lipschitz, that is for any $M>0$, there exists a constant $L_{M,m}<\infty$ such that for all $x,y\in[-M,M]^m$,
$
|f(\x)-f(\y)| \leq L_{M,m}\|\x-\y\|_2.
$
Further, $L_{M,m}\leq c(1+(M\sqrt{m})^{k-1})$ for some costant $c$.
\end{enumerate}

Using the above properties, we prove the following two technical lemmas used in the proof of Theorem \ref{thm:master_W2}

\begin{lemma}\label{lem:fL}
Let $g:\R\rightarrow\R$ be a Lipschitz function. Consider the function $f:\R^3\rightarrow\R$ defined as follows:
$$
f(\x) = x_1(x_2-g(x_3))^2.
$$
Then, $f\in\rm{PL}(3).$
\end{lemma}

\begin{proof}
Let $h:\R^2\rightarrow\R$ defined as $h(\ub)=(\ub_1-g(\ub_2))^2$. The function $(\ub_1,\ub_2)\mapsto\ub_1-g(\ub_2)$ is clearly Lipschitz. Thus, $h\in\rm{PL}(2)$, i.e., 
\begin{align}
|h(\ub)-h(\vb)| \leq C(1+\|\ub\|_2+\|\vb\|_2)\|\ub-\vb\|_2\quad\text{and}\quad |h(\vb)|\leq C'(1+\|\vb\|_2^2).\label{eq:h_pl}
\end{align} 
Therefore, letting $\x=(x_1,\ub)\in\R^3$ and $\y=(y_1,\vb)\in\R^3$, we have that 
\begin{align}
|f(\x)-f(\y)| &= |x_1h(\ub) - y_1h(\vb)| \leq |x_1||h(\ub)-h(\vb)| + |h(\vb)| |x_1-y_1|\nn\\
%&\leq |x_1|(1+\|\ub\|_2+\|\vb\|_2)\|\ub-\vb\|_2 + |h(\vb)| |x_1-y_1| \nn\\
&\leq C|x_1|(1+\|\ub\|_2+\|\vb\|_2)\|\ub-\vb\|_2 + C'(1+\|\vb\|_2^2)|x_1-y_1| \nn\\
&\leq C(|x_1|^2+(1+\|\ub\|_2+\|\vb\|_2)^2)\|\ub-\vb\|_2 + C'(1+\|\vb\|_2^2)|x_1-y_1| \nn\\
&\leq C(1+\|\x\|_2^2+\|\y\|_2^2)\|\ub-\vb\|_2 + C'(1+\|\x\|_2^2+\|\y\|_2^2)|x_1-y_1| \nn\\
&\leq C(1+\|\x\|_2^2+\|\y\|_2^2)\|\x-\y\|_2.
\end{align}
In the second line, we used \eqref{eq:h_pl}. In the third line, we used $2xy\leq x^2+y^2$. In the fourth line, we used Cauchy-Schwarz inequality. $C,C'>0$ are absolute constants that may change from line to line.

This completes the proof of the lemma.
\end{proof}

\fx{
\begin{lemma}[PL with Bounded Variables]\label{lem:PLbdd}%with bounded derivatives 
Let $f:\R^{d_1}\rightarrow \R$ be a $PL(k)$ function. $M\subset \R^{d_2}$ be a compact set and $g$ be a continuously differentiable function over $M$. Then $h(\x,\y)=f(\x)g(\y)$ is $PL(k+1)$ over $\R^{d_1}\times M$.
\end{lemma}
\begin{proof} First observe that since $g$ has continuous derivatives and is continuous over a compact set. Thus $g$ and its gradient is bounded and $g$ is Lipschitz over $M$. Let $B=\sup_{\x\in M}\max |g(x)|,\tn{\nabla g(\x)}$. To proceed, given pairs $(\x,\y)$ and $(\x',\y')$ over $\R\times M$, we have that
\begin{align}
|h(\x,\y)-h(\x',y')|&\leq |h(\x,\y)-h(\x',\y)|+ |h(\x',\y)-h(\x',\y')|\\
&\leq |f(\x)-f(\x')||g(\y)|+ |f(\x')||g(\y)-g(\y')|\\
&\leq B|f(\x)-f(\x')|+ B\tn{\y-\y'}|f(\x')|\\
&\leq B(1+\tn{\x'}^{k-1}+\tn{\x}^{k-1})\tn{\x-\x'}+ B\tn{\y-\y'}(1+\tn{\x'}^k)\\
&\lesssim (1+\tn{\z}^k+\tn{\z'}^{k})\tn{\z-\z'},
\end{align}
where $\z=[\x~\y]$. This shows the desired $\text{PL}(k+1)$ guarantee.
\end{proof}
The following lemma is in similar spirit to Lemma \ref{lem:PLbdd} and essentially follows from similar lines of arguments (i.e.~using Lipschitzness induced by boundedness).
}
\begin{lemma}\label{lem:hPL}
Let functions $f,g:\R^3\rightarrow\R$ such that $f\in\rm{PL}(k)$ and 
$$
g(x_1,x_2,x_3) := \frac{x_2^{-1/2}}{1+(\ksi_* x_2)^{-1}}\sqrt{\kappa}\tau_* x_3 + (1-(1+\ksi_*x_2)^{-1})x_1.
$$
Here, $\ksi_*,\tau_*,\kappa$ are positive constants. 
%bounded, say $x_2\in[m,M]\subset(0,\infty)$.
Further define 
\begin{align}
h(x,y,z) := f\left(g(y,z,x),y,z\right),
\end{align}
and assume that $y$ take values on a fixed bounded compact set $\Mc$.
Then, it holds that $h\in\rm{PL}(k+1)$.
\end{lemma}

\begin{proof}
Since $f$ is $\rm{PL}(k)$, for some $L>0$, \eqref{PL func} holds. Fix $x,x'\in\R$, $\ab=[y,z]\in\R^{2}$ and $\ab'=[y',z']\in\R^{2}$. Let $\bb=[\g,\ab]=[\g,y,z]\in\R^3$ where $\g=g(y,z,x)\in\R$ and define accordingly $\bb'$ and $\g'$. We have that
\begin{align}
|h([x,\ab])-h([x,\ab'])|&= |f(\bb)-f(\bb')|\nn\\
&\leq L\left(1+\|\bb\|_2^{k-1}+\|\bb'\|_2^{k-1}\right)\|\bb-\bb'\|_2\nn\\
&\leq C\left(1+\|\ab\|_2^{k-1}+\|\ab'\|_2^{k-1}+|\g|^{k-1}+|\g'|^{k-1}\right)(\|\ab-\ab'\|_2+|\g-\g'|),\label{eq:hlip}
%~~~~~~~~&L\left(1+\|\g(\y,\z,\x)\|_2^{k-1}+\|g(\y',\z',\x')\|_2^{k-1}\right)\|g(\y,\z,\x)-g(\y',\z',\x')\|_2.
\end{align}
for some constant $C>0$. In the last inequality we have repeatedly used the inequality
$
\left(\sum_{i=1}^m \|\vb_i\|_2^2\right)^{\frac{d}{2}} \leq  C(m)\cdot\sum_{i=1}^m\|\vb_i\|_2^{d}.
$
Next, we need to bound the $\g$ term in terms of $\ab$. This is accomplished as follows
\begin{align}
|\g|^{k-1} &= \left|\frac{y^{-1/2}}{1+(\ksi y)^{-1}}\sqrt{\kappa}\tau_* x + (1-(1+\ksi_*y)^{-1})z\right|^{k-1}\nn\\
%\|\g\|_2^{k-1}&\leq C \frac{1}{p}\sum_{i=1}^p|g_i|^{k-1}\\
%&\lesssim \frac{1}{p}\sum_{i=1}^p\left|\frac{y_i^{-1/2}}{1+(\ksi y_i)^{-1}}\sqrt{\kappa}\tau_* z_i + (1-(1+\ksi_*y_i)^{-1})x_i\right|\\
&\leq C (|z|+|x|)^{k-1} \nn\\
&\leq C \left(|x|^{k-1}+|z|^{k-1}\right) \leq C (|x|^{k-1} + \|\ab\|_2^{k-1}).\label{eq:gb}
\end{align}
Here, the value of the constant $C>0$ may change from line to line. 
Secondly and similarly, we have the following perturbation bound on the $\g-\g'$. Recall that $(x,y)$ and $(x',y')$ are bounded. We will prove the following sequence of inequalities
% \so{Suppose the triples $(x_i,y_i,z_i)$ takes values in a fixed bounded compact set $\Mc$. We will prove the following sequence of inequalities} 
\begin{align}
%|\|\g\|_2-\|\g'\|_2|&\leq 
|{\g-\g'}|&= |g(y,z,x)-g(y',z',x')| \nn\\
&\leq \left|\frac{y^{-1/2}}{1+(\ksi y)^{-1}}\sqrt{\kappa}\tau_* x - \frac{(y')^{-1/2}}{1+(\ksi y')^{-1}}\sqrt{\kappa}\tau_* x'\right|
 + \left|(1-(1+\ksi_*y)^{-1})z-(1-(1+\ksi_*y')^{-1})z'\right|\nn\\
 &\leq
\left|\frac{y^{-1/2}}{1+(\ksi y)^{-1}}\sqrt{\kappa}\tau_* x - \frac{y^{-1/2}}{1+(\ksi y)^{-1}}\sqrt{\kappa}\tau_* x'\right| +\left|\frac{y^{-1/2}}{1+(\ksi y)^{-1}}\sqrt{\kappa}\tau_* x'- \frac{(y')^{-1/2}}{1+(\ksi y')^{-1}}\sqrt{\kappa}\tau_* x'\right|\nn\\
 &\qquad+ \left|(1-(1+\ksi_*y)^{-1})z-(1-(1+\ksi_*y)^{-1})z'\right| + \left|(1-(1+\ksi_*y)^{-1})z'-(1-(1+\ksi_*y')^{-1})z'\right| \nn\\
&\leq C_1|x-x'| + C_2|x'||y-y'| + C_3|z-z'| + C_4|z'||y-y'|\nn\\
&\leq C(1+|x'| + |z'|)(|x-x'| + |z-z'| + |y-y'|)\\
&\leq C\sqrt{3}(1+|x'| + |z'|)\|[\ab,x]-[\ab',x']\|_2
.\label{eq:gd}
%C^2_{\Mc} \sum_{i=1}^p(|x_i-x'_i|^2+|y_i-y'_i|^2+|z_i-z'_i|^2)\label{sec line eq}\\
%&\leq C^2_{\Mc} (\tn{\x-\x'}^2+\tn{\y-\y'}^2+\tn{\z-\z'}^2)\\
%&\lesssim \tn{\ab-\ab'}^2.
\end{align}
In the fourth inequality above, we used the fact the assumption that $|y|$ is bounded. In the last line, we used Cauchy-Scwhartz.

Substituting \eqref{eq:gb} and \eqref{eq:gd} in \eqref{eq:hlip} gives:
\begin{align}
|h(x,y,z) - h(x',y',z')|  &\leq C\left(1+\|\ab\|_2^{k-1}+\|\ab'\|_2^{k-1}+|x|^{k-1}+|x'|^{k-1}\right)(1+|x'| + |z'|)\|[\ab,x]-[\ab',x']\|_2\nn \\
&\leq C\left(1+\|[\ab,x]\|_2^{k}+\|[\ab',x']\|_2^{k}\right)\|[\ab,x]-[\ab',x']\|_2.
\end{align}
Thus, $h\in\rm{PL}(k+1)$, as desired.
%In what follows, we prove the second line \eqref{sec line eq} i.e.~the fact that for any triples $a=(x,y,z),a'=(x',y',z')$ (with $a,a'\in\R^3$), we have that 
%\[
%|g(a)-g(a')|\leq C_{\Mc}\tn{a-a'}.
%\]
%Set $C_\nabla=\sup_{a\in \Mc}\tn{\nabla g(a)}$. By definition of gradient, the inequality above holds with $C_\Mc=C_\nabla$. Thus, all that remains is proving that $C_\nabla$ is upper bounded by a constant. \so{However, this automatically holds because from the definition of $g(\dots)$ function, it is clear that $C_\nabla$ is defined everywhere and continuous thus it has a finite maximum over a compact set.}
\end{proof}






















